动态规划和背包深度题卡 按题型拆成章内大节，但每个大节都先从 LeetCode 案例切入，再进入分流、案例矩阵、案例推导和实验复盘。这样读者看到的是从具体题推到规律，而不是先读抽象方法再猜它能用在哪。

\topic{动态规划与状态定义}
这一节先看 LeetCode 10 正则表达式匹配、LeetCode 44 通配符匹配、LeetCode 53 最大子数组和 这几道 LeetCode 题，再整理同一题型的分流和推导。阅读顺序不是先背原理，而是先看案例的暴力瓶颈，再把重复出现的观察抽象成方法。

\subsection{原理和适用条件}

以 LeetCode 10 正则表达式匹配、LeetCode 44 通配符匹配、LeetCode 53 最大子数组和 为主线：LeetCode 10 正则表达式匹配：模型是“模式中的点号匹配单字符，星号修饰前一个字符并允许出现零次或多次。”，优化抓手是“二维 DP 表示两个前缀是否匹配，星号转移同时覆盖零次和消耗一个字符后继续匹配。”；LeetCode 44 通配符匹配：模型是“问号匹配一个字符，星号匹配任意长度字符串，状态是两个前缀能否匹配。”，优化抓手是“DP 中星号可匹配空串或继续吞掉当前字符；贪心写法记录最近星号并在失配时回退扩展星号覆盖范围。”；LeetCode 53 最大子数组和：模型是“给定整数数组，返回一个非空连续子数组能够取得的最大元素和。”，优化抓手是“用 [-2, 3] 和 [2, 3] 对比，推出当前位置只在单独开始与接上旧后缀之间取较大值。”。这些案例共同推出的规律是：先写递归搜索并找重复状态，再给状态命名。状态必须包含未来决策所需信息，遍历顺序必须保证依赖状态已经算完。 方法名只是复盘后的标签，真正要掌握的是从题面约束、暴力失败、状态设计到边界反例的推导链。

\begin{principlebox}{图示：从 LeetCode 案例到 动态规划与状态定义推导链}
\begin{tabular}{@{}p{0.18\linewidth}p{0.74\linewidth}@{}}
暴力基线 & 枚举候选、完整模拟或逐个检查，先保证覆盖全部可能答案。\\
瓶颈定位 & 慢点不是递归写法，而是相同状态被反复计算。若一个状态的未来只取决于少量参数，就应保存结果并按依赖顺序计算。\\
结构选择 & 状态定义要包含足够信息但不能过度膨胀。线性 DP 看最后一步选择，二维 DP 看两个前缀或网格位置，背包看物品阶段和容量，区间 DP 看最后合并点。\\
不变量 & 填表不变量是：计算当前状态时，所有依赖状态已经算完且语义一致。滚动数组必须解释当前格读到的是上一轮还是本轮。\\
代码落地 & 数组维度和初始化要服务于状态含义。不可达用大正数或负数哨兵时要避免溢出。方案数可能超过 int 时改用 \code{long long} 或按题目取模。\\
\end{tabular}
\end{principlebox}

\subsection{LeetCode 案例分流}

\begin{principlebox}{从 LeetCode 案例分流：动态规划与状态定义}
\textbf{线性状态。}
代表案例：LeetCode 53 最大子数组和、LeetCode 70 爬楼梯、LeetCode 198 打家劫舍、LeetCode 300 LIS。先看这些题的题面条件：答案依赖前一个或少数几个位置的最优状态。由此推出的处理方式是：定义以当前位置结尾或考虑前 i 个对象的状态。风险点：状态含义不同，答案位置也不同。

\textbf{网格 DP。}
代表案例：LeetCode 62 不同路径、LeetCode 63 不同路径 II、LeetCode 64 最小路径和、LeetCode 221 最大正方形。先看这些题的题面条件：只能从固定方向进入格子。由此推出的处理方式是：状态表示到达当前格子的方案数或最优值。风险点：允许四方向移动时会变成图最短路。

\textbf{双字符串 DP。}
代表案例：LeetCode 72 编辑距离、LeetCode 97 交错字符串、LeetCode 1143 LCS、LeetCode 115 不同子序列。先看这些题的题面条件：两个前缀之间存在匹配、编辑或子序列关系。由此推出的处理方式是：dp[i][j] 表示两个前缀的答案。风险点：第零行第零列代表空串，不是随便补的边界。

\textbf{状态机 DP。}
代表案例：LeetCode 121 买卖股票 I、LeetCode 123 买卖股票 III、LeetCode 309 冷冻期、LeetCode 714 手续费。先看这些题的题面条件：每天或每步有持有、空闲、冷冻、交易次数等状态。由此推出的处理方式是：把每个状态看成日终状态，转移只能来自合法上一状态。风险点：同一天状态污染和交易次数维度是主要坑。

\textbf{区间和树形 DP。}
代表案例：LeetCode 312 戳气球、LeetCode 337 打家劫舍 III、LeetCode 124 最大路径和。先看这些题的题面条件：最后一步是区间内某个切分点，或父节点合并子树状态。由此推出的处理方式是：按区间长度或后序顺序计算。风险点：返回值和全局答案常常不是同一个量。

\end{principlebox}

\subsection{案例矩阵}

本节案例覆盖 LeetCode 10 正则表达式匹配、LeetCode 44 通配符匹配、LeetCode 53 最大子数组和、LeetCode 62 不同路径、LeetCode 63 不同路径 II、LeetCode 64 最小路径和、LeetCode 70 爬楼梯、LeetCode 72 编辑距离、LeetCode 97 交错字符串、LeetCode 115 不同的子序列、LeetCode 121 买卖股票的最佳时机、LeetCode 123 买卖股票 III、LeetCode 139 单词拆分、LeetCode 152 乘积最大子数组、LeetCode 188 买卖股票 IV、LeetCode 198 打家劫舍、LeetCode 221 最大正方形、LeetCode 300 最长递增子序列、LeetCode 309 最佳买卖股票含冷冻期、LeetCode 312 戳气球、LeetCode 329 矩阵中的最长递增路径、LeetCode 486 预测赢家、LeetCode 509 斐波那契数、LeetCode 516 最长回文子序列、LeetCode 600 不含连续 1 的非负整数、LeetCode 714 买卖股票含手续费、LeetCode 790 多米诺和托米诺平铺、LeetCode 1143 最长公共子序列、LeetCode 1235 规划兼职工作、LeetCode 1425 带限制的子序列和。阅读时不要按题号背答案，而要先判断题面落在哪个分流场景：查询目标是什么，输入约束给了什么单调性或等价关系，暴力慢在哪里，优化结构保存了什么状态。

\begin{principlebox}{案例矩阵}
\textbf{LeetCode 10 正则表达式匹配。}
模式中的点号匹配单字符，星号修饰前一个字符并允许出现零次或多次。单点风险：空串、能匹配空串的星号链、星号前必须有字符、点号和普通字符的匹配条件。

\textbf{LeetCode 44 通配符匹配。}
问号匹配一个字符，星号匹配任意长度字符串，状态是两个前缀能否匹配。单点风险：连续星号、空串、模式只含星号、贪心回退时字符串位置不能倒退错。

\textbf{LeetCode 53 最大子数组和。}
给定整数数组，返回一个非空连续子数组能够取得的最大元素和。单点风险：全负数组、单元素、和溢出、答案不能初始化为零。

\textbf{LeetCode 62 不同路径。}
网格状态由上方和左方两种最后一步转移而来。单点风险：一行、一列、较大网格结果溢出、障碍物版本边界。

\textbf{LeetCode 63 不同路径 II。}
障碍物让某些格子的路径数变为零。单点风险：起点或终点是障碍、整行被截断、单行障碍。

\textbf{LeetCode 64 最小路径和。}
状态是到达当前格子的最小代价，最后一步来自上方或左方。单点风险：单行、单列、权值为零、路径和溢出。

\textbf{LeetCode 70 爬楼梯。}
最后一步来自前一阶或前两阶，是最小的重叠子问题模型。单点风险：n 为一、n 为二、题目是否允许零阶、结果范围。

\textbf{LeetCode 72 编辑距离。}
二维状态表示两个前缀之间的最少编辑操作。单点风险：空串、完全相同、完全不同、操作代价不同。

\textbf{LeetCode 97 交错字符串。}
目标串前缀由两个来源字符串的前缀交错组成。单点风险：总长度不等、第一行第一列初始化、两个来源字符都能匹配时不能贪心只选一个。

\textbf{LeetCode 115 不同的子序列。}
从源串中删除若干字符后形成目标串，本质是两个前缀之间的计数问题。单点风险：空目标串计数为一、源串为空、重复字符导致计数增长、计数可能超过 int。

\textbf{LeetCode 121 买卖股票的最佳时机。}
卖出日固定时，最佳买入日是此前最低价格。单点风险：单日价格、一直下跌、一直上涨、价格相同。

\textbf{LeetCode 123 买卖股票 III。}
最多两次交易，状态需要包含交易次数和持有状态。单点风险：价格为空、只够一次交易、两次交易间不能重叠。

\textbf{LeetCode 139 单词拆分。}
状态表示前缀能否由字典单词拼出。单点风险：空串、重复单词、长单词、substr 复制成本。

\textbf{LeetCode 152 乘积最大子数组。}
负数会让最大乘积和最小乘积互换。单点风险：零、负数个数奇偶、单元素、乘积溢出。

\textbf{LeetCode 188 买卖股票 IV。}
最多 k 次交易让状态必须同时包含交易次数和是否持有股票。单点风险：k 为零、价格天数很少、k 大于等于 n/2、更新顺序造成同一天状态污染。

\textbf{LeetCode 198 打家劫舍。}
每间房只有偷和不偷两种结局，偷当前就不能偷前一间。单点风险：空数组、单间、全零、金额很大。

\textbf{LeetCode 221 最大正方形。}
以当前位置为右下角的最大正方形由上、左、左上三个状态限制。单点风险：全零、全一、单行单列、字符一和数字一不要混淆。

\textbf{LeetCode 300 最长递增子序列。}
二次 DP 固定结尾，二分优化维护每个长度的最小结尾。单点风险：重复元素、严格递增和非递减区别、空数组。

\textbf{LeetCode 309 最佳买卖股票含冷冻期。}
冷冻期让买入来源受昨天卖出状态限制。单点风险：空价格、一天、连续上涨、卖出后不能次日买入。

\textbf{LeetCode 312 戳气球。}
先戳哪个会让邻居持续变化，改成最后戳某个气球后区间边界才稳定。单点风险：补一的虚拟边界、区间长度遍历顺序、空区间收益为零、乘积可能较大。

\textbf{LeetCode 329 矩阵中的最长递增路径。}
每个格子的答案是从它出发能走出的最长严格递增路径。单点风险：平台相等不能走、四方向越界、递归深度、允许不下降时会产生环。

\textbf{LeetCode 486 预测赢家。}
两人从区间两端取数，当前选择会影响对手面对的剩余区间。单点风险：单元素、两个元素、总分相等、负数若存在、区间遍历顺序。

\textbf{LeetCode 509 斐波那契数。}
递归定义会让 fib(n-1) 和 fib(n-2) 下面的状态反复计算。单点风险：n 为零、n 为一、结果范围、递归爆栈。

\textbf{LeetCode 516 最长回文子序列。}
子序列允许跳过字符，区间两端是否相等决定能否同时作为回文两端。单点风险：空串、单字符、两端相等但内部为空、全相同字符、子串和子序列混淆。

\textbf{LeetCode 600 不含连续 1 的非负整数。}
统计不超过 n 的二进制数，需要按位构造并记住上一位是否为一。单点风险：n 为零、最高位、连续一提前停止、包含 n 本身、整数位宽。

\textbf{LeetCode 714 买卖股票含手续费。}
手续费改变交易收益，仍可用持有和不持有状态。单点风险：手续费为零、价格震荡、小利润不足手续费。

\textbf{LeetCode 790 多米诺和托米诺平铺。}
铺满前 i 列时，边界可能完整，也可能上缺或下缺一个格子。单点风险：n 为一、n 为二、取模、缺口上下对称、初始化。

\textbf{LeetCode 1143 最长公共子序列。}
状态表示两个前缀的最长公共子序列长度。单点风险：空串、完全相同、无公共字符、子序列不是子串。

\textbf{LeetCode 1235 规划兼职工作。}
每份工作有开始、结束和收益，选择不冲突工作使收益最大。单点风险：结束等于开始是否兼容、收益相同、工作排序、空列表。

\textbf{LeetCode 1425 带限制的子序列和。}
dp[i] 只依赖前 k 个位置的最大 dp 值，转移范围是滑动窗口。单点风险：全负数、k 为一、k 大于数组长度、答案必须非空、队列过期顺序。

\end{principlebox}

\subsection{共性落点}

\begin{principlebox}{本组共性落点}
\textbf{慢版本。} 暴力搜索枚举选择序列、切分点或路径，很多不同路径会到达同一个子问题。DP 的第一步是给这些重复子问题命名。
\textbf{瓶颈。} 慢点不是递归写法，而是相同状态被反复计算。若一个状态的未来只取决于少量参数，就应保存结果并按依赖顺序计算。
\textbf{状态字段。} 状态下标、状态值语义、初始化、转移来源、遍历顺序；空间压缩时还要指出读到的是旧值还是新值。
\textbf{不变量。} 填表不变量是：计算当前状态时，所有依赖状态已经算完且语义一致。滚动数组必须解释当前格读到的是上一轮还是本轮。
\textbf{实现检查。} 先写未压缩版本校验转移；压缩前后用小表对拍；不可达哨兵参与加法前要检查溢出。
\textbf{C++ 落地。} 数组维度和初始化要服务于状态含义。不可达用大正数或负数哨兵时要避免溢出。方案数可能超过 int 时改用 \code{long long} 或按题目取模。
\textbf{证明抓手。} 证明从最后一步开始：任意最优解的最后一步都在转移枚举中，转移来源已经由更小状态正确计算。
\end{principlebox}

\subsection{案例推导}

\noindent\textbf{LeetCode 10 正则表达式匹配。}

\textbf{题目说明。} LeetCode 10 正则表达式匹配 的题面入口要先还原成具体任务：模式中的点号匹配单字符，星号修饰前一个字符并允许出现零次或多次。

\textbf{样例入口。} 拿 s=a、p=a* 手算，星号可以让 a 出现一次；拿 s=空、p=a*，星号又可以让 a 出现零次。再拿 s=aa、p=a*，消耗一个 a 后模式仍停在 a*，还能继续匹配。

\textbf{暴力基线。} 慢版本先写递归搜索树：节点表示处理位置、剩余资源和当前选择。只要不同路径反复到达同一个节点，就说明需要把这个节点命名成 DP 状态。放到本题就是：先按“模式中的点号匹配单字符，星号修饰前一个字符并允许出现零次或多次”完整试慢版本；样例里已经能看到“规律是星号有两条分叉：当零次用时看 dp[i][j-2]，当消耗当前字符时看 dp[i-1][j]”，所以优化前必须先能说清慢版本枚举了哪些候选。

\textbf{暴力过程。} 拿 s=a、p=a* 手算，星号可以让 a 出现一次；拿 s=空、p=a*，星号又可以让 a 出现零次。再拿 s=aa、p=a*，消耗一个 a 后模式仍停在 a*，还能继续匹配。

\textbf{重复工作。} 题面模型：模式中的点号匹配单字符，星号修饰前一个字符并允许出现零次或多次。暴力版的慢点要落到本题：慢版本会在“模式中的点号匹配单字符，星号修饰前一个字符并允许出现零次或多次”的选择树里反复到达相同参数。样例规律是“规律是星号有两条分叉：当零次用时看 dp[i][j-2]，当消耗当前字符时看 dp[i-1][j]”，说明这个参数组合可以命名成状态，算过之后后面直接复用。后面的优化只消掉这类重复工作，不能改变答案集合。

\textbf{优化条件。} 本题的关键观察是：二维 DP 表示两个前缀是否匹配，星号转移同时覆盖零次和消耗一个字符后继续匹配。这句话必须能翻译成可维护状态；如果题目条件改变后观察不再成立，就不能继续套原来的写法。

\textbf{相邻状态。} 规律是星号有两条分叉：当零次用时看 dp[i][j-2]，当消耗当前字符时看 dp[i-1][j]。把这个特例继续拆开看：先问暴力搜索里哪些不同路径到达同一个后续问题，再给这个后续问题命名为状态。状态必须包含未来决策需要的全部信息，转移只从更小且已经算清的状态来。

\textbf{状态复用。} 把重复递归节点命名为状态；遍历顺序只是在保证依赖状态已经算完。本题落点是：规律是星号有两条分叉：当零次用时看 dp[i][j-2]，当消耗当前字符时看 dp[i-1][j]。手算时只改被新输入影响的那一小部分，而不是回头重算完整候选。

\textbf{指针和字段。} 状态字段要能说清：状态下标、状态值语义、初始化、转移来源、遍历顺序；空间压缩时还要指出读到的是旧值还是新值。手算时按这个协议跟踪：拿 s=a、p=a* 手算，星号可以让 a 出现一次；拿 s=空、p=a*，星号又可以让 a 出现零次。再拿 s=aa、p=a*，消耗一个 a 后模式仍停在 a*，还能继续匹配。然后按协议复查：手算时先写一个很小的 DP 表或记忆化搜索记录。每个格子旁边写它来自哪些更小状态。

\textbf{代码落点。} 先写未压缩版本校验转移；压缩前后用小表对拍；不可达哨兵参与加法前要检查溢出。关键代码行必须能逐句翻译成“旧状态怎样变成新状态”，否则就只是模板。

\textbf{正确性。} 证明从最后一步开始：任意最优解的最后一步都在转移枚举中，转移来源已经由更小状态正确计算。

\textbf{边界实验。} 先用空串、能匹配空串的星号链、星号前必须有字符、点号和普通字符的匹配条件做反例检查。实验时覆盖空串、能匹配空串的星号链、星号前必须有字符、点号和普通字符的匹配条件，先打印完整 DP 表，再比较空间压缩版本。若压缩版不同，优先检查遍历顺序和旧值保存。

\textbf{复杂度变化。} 暴力递归会重复计算相同子问题，常见指数级；DP 把每个状态算一次，时间是状态数乘单个转移成本，空间是状态表大小或压缩后的大小。

\textbf{迁移追问。} 如果题目改成“通配符匹配中的星号语义是任意长度字符串，不能直接搬用本题转移”，先判断原状态是否仍然覆盖新答案集合，再决定是微调边界、改 value 语义，还是换算法。

\noindent\textbf{LeetCode 44 通配符匹配。}

\textbf{题目说明。} LeetCode 44 通配符匹配 的题面入口要先还原成具体任务：问号匹配一个字符，星号匹配任意长度字符串，状态是两个前缀能否匹配。

\textbf{样例入口。} 拿 s=ab、p=a* 手算，星号可以吞掉 b；拿 s=空、p=*，星号也可以代表空。DP 中星号的两条路是匹配空串或继续吞当前字符。贪心中失配时回到最近星号，让它多吞一个字符。

\textbf{暴力基线。} 慢版本先写递归搜索树：节点表示处理位置、剩余资源和当前选择。只要不同路径反复到达同一个节点，就说明需要把这个节点命名成 DP 状态。放到本题就是：先按“问号匹配一个字符，星号匹配任意长度字符串，状态是两个前缀能否匹配”完整试慢版本；样例里已经能看到“规律是本题星号独立代表任意串，和正则里修饰前一字符的星号不同”，所以优化前必须先能说清慢版本枚举了哪些候选。

\textbf{暴力过程。} 拿 s=ab、p=a* 手算，星号可以吞掉 b；拿 s=空、p=*，星号也可以代表空。DP 中星号的两条路是匹配空串或继续吞当前字符。贪心中失配时回到最近星号，让它多吞一个字符。

\textbf{重复工作。} 题面模型：问号匹配一个字符，星号匹配任意长度字符串，状态是两个前缀能否匹配。暴力版的慢点要落到本题：慢版本会在“问号匹配一个字符，星号匹配任意长度字符串，状态是两个前缀能否匹配”的选择树里反复到达相同参数。样例规律是“规律是本题星号独立代表任意串，和正则里修饰前一字符的星号不同”，说明这个参数组合可以命名成状态，算过之后后面直接复用。后面的优化只消掉这类重复工作，不能改变答案集合。

\textbf{优化条件。} 本题的关键观察是：DP 中星号可匹配空串或继续吞掉当前字符；贪心写法记录最近星号并在失配时回退扩展星号覆盖范围。这句话必须能翻译成可维护状态；如果题目条件改变后观察不再成立，就不能继续套原来的写法。

\textbf{相邻状态。} 规律是本题星号独立代表任意串，和正则里修饰前一字符的星号不同。把这个特例继续拆开看：先问暴力搜索里哪些不同路径到达同一个后续问题，再给这个后续问题命名为状态。状态必须包含未来决策需要的全部信息，转移只从更小且已经算清的状态来。

\textbf{状态复用。} 把重复递归节点命名为状态；遍历顺序只是在保证依赖状态已经算完。本题落点是：规律是本题星号独立代表任意串，和正则里修饰前一字符的星号不同。手算时只改被新输入影响的那一小部分，而不是回头重算完整候选。

\textbf{指针和字段。} 状态字段要能说清：状态下标、状态值语义、初始化、转移来源、遍历顺序；空间压缩时还要指出读到的是旧值还是新值。手算时按这个协议跟踪：拿 s=ab、p=a* 手算，星号可以吞掉 b；拿 s=空、p=*，星号也可以代表空。DP 中星号的两条路是匹配空串或继续吞当前字符。贪心中失配时回到最近星号，让它多吞一个字符。然后按协议复查：手算时先写一个很小的 DP 表或记忆化搜索记录。每个格子旁边写它来自哪些更小状态。

\textbf{代码落点。} 先写未压缩版本校验转移；压缩前后用小表对拍；不可达哨兵参与加法前要检查溢出。关键代码行必须能逐句翻译成“旧状态怎样变成新状态”，否则就只是模板。

\textbf{正确性。} 证明从最后一步开始：任意最优解的最后一步都在转移枚举中，转移来源已经由更小状态正确计算。

\textbf{边界实验。} 先用连续星号、空串、模式只含星号、贪心回退时字符串位置不能倒退错做反例检查。实验时覆盖连续星号、空串、模式只含星号、贪心回退时字符串位置不能倒退错，先打印完整 DP 表，再比较空间压缩版本。若压缩版不同，优先检查遍历顺序和旧值保存。

\textbf{复杂度变化。} 暴力递归会重复计算相同子问题，常见指数级；DP 把每个状态算一次，时间是状态数乘单个转移成本，空间是状态表大小或压缩后的大小。

\textbf{迁移追问。} 如果题目改成“正则表达式匹配的星号修饰前一个字符，和本题星号独立匹配任意串不同”，先判断原状态是否仍然覆盖新答案集合，再决定是微调边界、改 value 语义，还是换算法。

\noindent\textbf{LeetCode 53 最大子数组和。}

\textbf{题目说明。} LeetCode 53 最大子数组和给定整数数组 nums，要求找一个非空连续子数组，使元素和最大，并返回这个最大和。

\textbf{样例入口。} 样例 nums=[-2, 1, -3, 4, -1, 2, 1, -5, 4]，输出 6；连续段 [4, -1, 2, 1] 的和为 6。

\textbf{暴力基线。} 暴力枚举所有连续子数组，累计每段和并更新最大值；优化一点可以固定左端、右端扩展时维护当前和。

\textbf{暴力过程。} 以 [-2, 1, -3, 4, -1, 2, 1] 为例，暴力会重复计算很多以 4 开头或以 2 结尾的子段。真正关键是当前数要不要接上前一个后缀。

\textbf{重复工作。} 题面模型：给定整数数组，返回一个非空连续子数组能够取得的最大元素和。暴力版的慢点要落到本题：浪费在重复求每个后缀和。对于以 i 结尾的最优连续段，历史信息只需要一个值：以 i-1 结尾的最大后缀和。后面的优化只消掉这类重复工作，不能改变答案集合。

\textbf{优化条件。} 本题的关键观察是：用 [-2, 3] 和 [2, 3] 对比，推出当前位置只在单独开始与接上旧后缀之间取较大值。这句话必须能翻译成可维护状态；如果题目条件改变后观察不再成立，就不能继续套原来的写法。

\textbf{相邻状态。} 规律是当前位置只需要比较 nums[i] 与 bestEndingHere + nums[i]，全局答案再和 bestEndingHere 比较；全负数组要求答案初始化为第一个数，不能初始化为 0。把这个特例继续拆开看：先问暴力搜索里哪些不同路径到达同一个后续问题，再给这个后续问题命名为状态。状态必须包含未来决策需要的全部信息，转移只从更小且已经算清的状态来。

\textbf{状态复用。} bestEndingHere 表示以当前位置结尾的最大子数组和；answer 表示全局最大值。当前位置要么单独开始，要么接上旧后缀。手算时只改被新输入影响的那一小部分，而不是回头重算完整候选。

\textbf{指针和字段。} 状态字段要能说清：current=bestEndingHere，answer 是全局最优；nums[i] 与 current+nums[i] 比较决定是否丢弃旧后缀。手算时按这个协议跟踪：遇到 4 之前，旧后缀为负，接上反而变差，所以从 4 重新开始；之后 -1、2、1 继续接上得到 6。

\textbf{代码落点。} 关键代码是 current=max(nums[i], current+nums[i]); answer=max(answer, current)。answer 必须初始化为 nums[0]，否则全负数组会错。关键代码行必须能逐句翻译成“旧状态怎样变成新状态”，否则就只是模板。

\textbf{正确性。} 任意以 i 结尾的连续段只有两类：只含 nums[i]，或由某个以 i-1 结尾的连续段加上 nums[i]。取前一类最优即可覆盖全部候选。

\textbf{边界实验。} 先用全负数组、单元素、和溢出、答案不能初始化为零做反例检查。实验时覆盖全负数组、单元素、和溢出、答案不能初始化为零，先打印完整 DP 表，再比较空间压缩版本。若压缩版不同，优先检查遍历顺序和旧值保存。

\textbf{复杂度变化。} 暴力 O(\ensuremath{n^{2}}) 或 O(\ensuremath{n^{3}})，Kadane 算法 O(n) 时间、O(1) 空间。

\textbf{迁移追问。} 如果题目改成“若要求返回区间下标，同步记录重新开始位置和最优左右端”，先判断原状态是否仍然覆盖新答案集合，再决定是微调边界、改 value 语义，还是换算法。

\noindent\textbf{LeetCode 62 不同路径。}

\textbf{题目说明。} LeetCode 62 不同路径 的题面入口要先还原成具体任务：网格状态由上方和左方两种最后一步转移而来。

\textbf{样例入口。} 拿 2x3 网格手算，到每个格子的最后一步只能从上方或左方来；第一行和第一列都只有 1 条路，右下角等于左边 2 加上上边 1，得到 3。

\textbf{暴力基线。} 慢版本先写递归搜索树：节点表示处理位置、剩余资源和当前选择。只要不同路径反复到达同一个节点，就说明需要把这个节点命名成 DP 状态。放到本题就是：先按“网格状态由上方和左方两种最后一步转移而来”完整试慢版本；样例里已经能看到“规律是 dp[row][col]=上方+左方，滚动数组只是在保存上一行同列和当前行左侧”，所以优化前必须先能说清慢版本枚举了哪些候选。

\textbf{暴力过程。} 拿 2x3 网格手算，到每个格子的最后一步只能从上方或左方来；第一行和第一列都只有 1 条路，右下角等于左边 2 加上上边 1，得到 3。

\textbf{重复工作。} 题面模型：网格状态由上方和左方两种最后一步转移而来。暴力版的慢点要落到本题：慢版本会在“网格状态由上方和左方两种最后一步转移而来”的选择树里反复到达相同参数。样例规律是“规律是 dp[row][col]=上方+左方，滚动数组只是在保存上一行同列和当前行左侧”，说明这个参数组合可以命名成状态，算过之后后面直接复用。后面的优化只消掉这类重复工作，不能改变答案集合。

\textbf{优化条件。} 本题的关键观察是：二维表可压缩为一维，因为当前行只依赖上一行同列和当前行左侧。这句话必须能翻译成可维护状态；如果题目条件改变后观察不再成立，就不能继续套原来的写法。

\textbf{相邻状态。} 规律是 dp[row][col]=上方+左方，滚动数组只是在保存上一行同列和当前行左侧。把这个特例继续拆开看：先问暴力搜索里哪些不同路径到达同一个后续问题，再给这个后续问题命名为状态。状态必须包含未来决策需要的全部信息，转移只从更小且已经算清的状态来。

\textbf{状态复用。} 把重复递归节点命名为状态；遍历顺序只是在保证依赖状态已经算完。本题落点是：规律是 dp[row][col]=上方+左方，滚动数组只是在保存上一行同列和当前行左侧。手算时只改被新输入影响的那一小部分，而不是回头重算完整候选。

\textbf{指针和字段。} 状态字段要能说清：状态下标、状态值语义、初始化、转移来源、遍历顺序；空间压缩时还要指出读到的是旧值还是新值。手算时按这个协议跟踪：拿 2x3 网格手算，到每个格子的最后一步只能从上方或左方来；第一行和第一列都只有 1 条路，右下角等于左边 2 加上上边 1，得到 3。然后按协议复查：手算时先写一个很小的 DP 表或记忆化搜索记录。每个格子旁边写它来自哪些更小状态。

\textbf{代码落点。} 先写未压缩版本校验转移；压缩前后用小表对拍；不可达哨兵参与加法前要检查溢出。关键代码行必须能逐句翻译成“旧状态怎样变成新状态”，否则就只是模板。

\textbf{正确性。} 证明从最后一步开始：任意最优解的最后一步都在转移枚举中，转移来源已经由更小状态正确计算。

\textbf{边界实验。} 先用一行、一列、较大网格结果溢出、障碍物版本边界做反例检查。实验时覆盖一行、一列、较大网格结果溢出、障碍物版本边界，先打印完整 DP 表，再比较空间压缩版本。若压缩版不同，优先检查遍历顺序和旧值保存。

\textbf{复杂度变化。} 暴力递归会重复计算相同子问题，常见指数级；DP 把每个状态算一次，时间是状态数乘单个转移成本，空间是状态表大小或压缩后的大小。

\textbf{迁移追问。} 如果题目改成“带障碍物时遇到障碍状态清零，第一行第一列要按阻断处理”，先判断原状态是否仍然覆盖新答案集合，再决定是微调边界、改 value 语义，还是换算法。

\noindent\textbf{LeetCode 63 不同路径 II。}

\textbf{题目说明。} LeetCode 63 不同路径 II 的题面入口要先还原成具体任务：障碍物让某些格子的路径数变为零。

\textbf{样例入口。} 拿 [[0, 0], [1, 0]] 手算，左下角是障碍，路径数必须清零；右下角只能从上方到达，所以答案是 1。

\textbf{暴力基线。} 慢版本先写递归搜索树：节点表示处理位置、剩余资源和当前选择。只要不同路径反复到达同一个节点，就说明需要把这个节点命名成 DP 状态。放到本题就是：先按“障碍物让某些格子的路径数变为零”完整试慢版本；样例里已经能看到“规律是障碍格不是普通转移，而是把当前状态置零；起点或终点是障碍时也按同一清零语义处理”，所以优化前必须先能说清慢版本枚举了哪些候选。

\textbf{暴力过程。} 拿 [[0, 0], [1, 0]] 手算，左下角是障碍，路径数必须清零；右下角只能从上方到达，所以答案是 1。

\textbf{重复工作。} 题面模型：障碍物让某些格子的路径数变为零。暴力版的慢点要落到本题：慢版本会在“障碍物让某些格子的路径数变为零”的选择树里反复到达相同参数。样例规律是“规律是障碍格不是普通转移，而是把当前状态置零；起点或终点是障碍时也按同一清零语义处理”，说明这个参数组合可以命名成状态，算过之后后面直接复用。后面的优化只消掉这类重复工作，不能改变答案集合。

\textbf{优化条件。} 本题的关键观察是：滚动数组中遇到障碍直接把当前位置清零，否则累加左侧。这句话必须能翻译成可维护状态；如果题目条件改变后观察不再成立，就不能继续套原来的写法。

\textbf{相邻状态。} 规律是障碍格不是普通转移，而是把当前状态置零；起点或终点是障碍时也按同一清零语义处理。把这个特例继续拆开看：先问暴力搜索里哪些不同路径到达同一个后续问题，再给这个后续问题命名为状态。状态必须包含未来决策需要的全部信息，转移只从更小且已经算清的状态来。

\textbf{状态复用。} 把重复递归节点命名为状态；遍历顺序只是在保证依赖状态已经算完。本题落点是：规律是障碍格不是普通转移，而是把当前状态置零；起点或终点是障碍时也按同一清零语义处理。手算时只改被新输入影响的那一小部分，而不是回头重算完整候选。

\textbf{指针和字段。} 状态字段要能说清：状态下标、状态值语义、初始化、转移来源、遍历顺序；空间压缩时还要指出读到的是旧值还是新值。手算时按这个协议跟踪：拿 [[0, 0], [1, 0]] 手算，左下角是障碍，路径数必须清零；右下角只能从上方到达，所以答案是 1。然后按协议复查：手算时先写一个很小的 DP 表或记忆化搜索记录。每个格子旁边写它来自哪些更小状态。

\textbf{代码落点。} 先写未压缩版本校验转移；压缩前后用小表对拍；不可达哨兵参与加法前要检查溢出。关键代码行必须能逐句翻译成“旧状态怎样变成新状态”，否则就只是模板。

\textbf{正确性。} 证明从最后一步开始：任意最优解的最后一步都在转移枚举中，转移来源已经由更小状态正确计算。

\textbf{边界实验。} 先用起点或终点是障碍、整行被截断、单行障碍做反例检查。实验时覆盖起点或终点是障碍、整行被截断、单行障碍，先打印完整 DP 表，再比较空间压缩版本。若压缩版不同，优先检查遍历顺序和旧值保存。

\textbf{复杂度变化。} 暴力递归会重复计算相同子问题，常见指数级；DP 把每个状态算一次，时间是状态数乘单个转移成本，空间是状态表大小或压缩后的大小。

\textbf{迁移追问。} 如果题目改成“若允许向上或向左走，图会有环，普通填表不再成立”，先判断原状态是否仍然覆盖新答案集合，再决定是微调边界、改 value 语义，还是换算法。

\noindent\textbf{LeetCode 64 最小路径和。}

\textbf{题目说明。} LeetCode 64 最小路径和 的题面入口要先还原成具体任务：状态是到达当前格子的最小代价，最后一步来自上方或左方。

\textbf{样例入口。} 拿 [[1, 3, 1], [1, 5, 1], [4, 2, 1]] 手算，到 5 的最小代价来自左侧 4 或上方 4，再加 5；右下角只需要比较上方和左方的已知最小代价。

\textbf{暴力基线。} 慢版本先写递归搜索树：节点表示处理位置、剩余资源和当前选择。只要不同路径反复到达同一个节点，就说明需要把这个节点命名成 DP 状态。放到本题就是：先按“状态是到达当前格子的最小代价，最后一步来自上方或左方”完整试慢版本；样例里已经能看到“规律是最后一步来自上或左，先初始化第一行第一列，内部统一取 min”，所以优化前必须先能说清慢版本枚举了哪些候选。

\textbf{暴力过程。} 拿 [[1, 3, 1], [1, 5, 1], [4, 2, 1]] 手算，到 5 的最小代价来自左侧 4 或上方 4，再加 5；右下角只需要比较上方和左方的已知最小代价。

\textbf{重复工作。} 题面模型：状态是到达当前格子的最小代价，最后一步来自上方或左方。暴力版的慢点要落到本题：慢版本会在“状态是到达当前格子的最小代价，最后一步来自上方或左方”的选择树里反复到达相同参数。样例规律是“规律是最后一步来自上或左，先初始化第一行第一列，内部统一取 min”，说明这个参数组合可以命名成状态，算过之后后面直接复用。后面的优化只消掉这类重复工作，不能改变答案集合。

\textbf{优化条件。} 本题的关键观察是：先初始化第一行和第一列，再处理内部，主转移保持简单。这句话必须能翻译成可维护状态；如果题目条件改变后观察不再成立，就不能继续套原来的写法。

\textbf{相邻状态。} 规律是最后一步来自上或左，先初始化第一行第一列，内部统一取 min。把这个特例继续拆开看：先问暴力搜索里哪些不同路径到达同一个后续问题，再给这个后续问题命名为状态。状态必须包含未来决策需要的全部信息，转移只从更小且已经算清的状态来。

\textbf{状态复用。} 把重复递归节点命名为状态；遍历顺序只是在保证依赖状态已经算完。本题落点是：规律是最后一步来自上或左，先初始化第一行第一列，内部统一取 min。手算时只改被新输入影响的那一小部分，而不是回头重算完整候选。

\textbf{指针和字段。} 状态字段要能说清：状态下标、状态值语义、初始化、转移来源、遍历顺序；空间压缩时还要指出读到的是旧值还是新值。手算时按这个协议跟踪：拿 [[1, 3, 1], [1, 5, 1], [4, 2, 1]] 手算，到 5 的最小代价来自左侧 4 或上方 4，再加 5；右下角只需要比较上方和左方的已知最小代价。然后按协议复查：手算时先写一个很小的 DP 表或记忆化搜索记录。每个格子旁边写它来自哪些更小状态。

\textbf{代码落点。} 先写未压缩版本校验转移；压缩前后用小表对拍；不可达哨兵参与加法前要检查溢出。关键代码行必须能逐句翻译成“旧状态怎样变成新状态”，否则就只是模板。

\textbf{正确性。} 证明从最后一步开始：任意最优解的最后一步都在转移枚举中，转移来源已经由更小状态正确计算。

\textbf{边界实验。} 先用单行、单列、权值为零、路径和溢出做反例检查。实验时覆盖单行、单列、权值为零、路径和溢出，先打印完整 DP 表，再比较空间压缩版本。若压缩版不同，优先检查遍历顺序和旧值保存。

\textbf{复杂度变化。} 暴力递归会重复计算相同子问题，常见指数级；DP 把每个状态算一次，时间是状态数乘单个转移成本，空间是状态表大小或压缩后的大小。

\textbf{迁移追问。} 如果题目改成“若允许四方向移动，变成最短路而不是普通网格 DP”，先判断原状态是否仍然覆盖新答案集合，再决定是微调边界、改 value 语义，还是换算法。

\noindent\textbf{LeetCode 70 爬楼梯。}

\textbf{题目说明。} LeetCode 70 爬楼梯 的题面入口要先还原成具体任务：最后一步来自前一阶或前两阶，是最小的重叠子问题模型。

\textbf{样例入口。} 拿 n=4 手算，最后一步要么从第 3 阶跨 1 阶，要么从第 2 阶跨 2 阶，这两类互斥且覆盖所有走法。于是 ways[4]=ways[3]+ways[2]。

\textbf{暴力基线。} 暴力递归枚举每一步走 1 阶或 2 阶，直到正好到 n 计数，超过 n 则失败。

\textbf{暴力过程。} 以 n=4 为例，最后一步要么从第 3 阶走 1 阶，要么从第 2 阶走 2 阶；这两类互斥且覆盖所有走法。

\textbf{重复工作。} 题面模型：最后一步来自前一阶或前两阶，是最小的重叠子问题模型。暴力版的慢点要落到本题：浪费在递归树反复计算 ways(2)、ways(3) 等同一子问题。到达某一阶后，后面的走法数量和之前怎么来无关。后面的优化只消掉这类重复工作，不能改变答案集合。

\textbf{优化条件。} 本题的关键观察是：滚动变量保存最近两个状态即可，不需要完整数组。这句话必须能翻译成可维护状态；如果题目条件改变后观察不再成立，就不能继续套原来的写法。

\textbf{相邻状态。} 规律不是背斐波那契，而是最后一步只有两个来源，状态只需保存前两项。把这个特例继续拆开看：先问暴力搜索里哪些不同路径到达同一个后续问题，再给这个后续问题命名为状态。状态必须包含未来决策需要的全部信息，转移只从更小且已经算清的状态来。

\textbf{状态复用。} dp[i] 表示到达第 i 阶的方法数，转移是 dp[i]=dp[i-1]+dp[i-2]。再压缩成两个滚动变量。手算时只改被新输入影响的那一小部分，而不是回头重算完整候选。

\textbf{指针和字段。} 状态字段要能说清：prev2 表示 dp[i-2]，prev1 表示 dp[i-1]，current 表示 dp[i]。手算时按这个协议跟踪：n=4 时 dp[1]=1、dp[2]=2、dp[3]=3、dp[4]=5。每一步只依赖前两项。

\textbf{代码落点。} 关键是初始化：n=0 按题意通常返回 1 或不出现；LeetCode 输入 n>=1，n=1 返回 1，n=2 返回 2。关键代码行必须能逐句翻译成“旧状态怎样变成新状态”，否则就只是模板。

\textbf{正确性。} 所有到第 i 阶的方案按最后一步分成从 i-1 来和从 i-2 来两类，二者不重不漏。

\textbf{边界实验。} 先用n 为一、n 为二、题目是否允许零阶、结果范围做反例检查。实验时覆盖n 为一、n 为二、题目是否允许零阶、结果范围，先打印完整 DP 表，再比较空间压缩版本。若压缩版不同，优先检查遍历顺序和旧值保存。

\textbf{复杂度变化。} 暴力递归指数级；DP O(n) 时间，滚动变量 O(1) 空间。

\textbf{迁移追问。} 如果题目改成“若每次可走一到 k 阶，转移变成固定宽度窗口求和”，先判断原状态是否仍然覆盖新答案集合，再决定是微调边界、改 value 语义，还是换算法。

\noindent\textbf{LeetCode 72 编辑距离。}

\textbf{题目说明。} LeetCode 72 编辑距离 的题面入口要先还原成具体任务：二维状态表示两个前缀之间的最少编辑操作。

\textbf{样例入口。} 拿 word1=ab、word2=ac 手算最后一个字符。若 b 改成 c，是 dp[1][1]+1；若删掉 b，是 dp[1][2]+1；若插入 c，是 dp[2][1]+1。字符相等时才直接继承左上。

\textbf{暴力基线。} 暴力递归比较两个字符串后缀：相等就同时跳过，不等就尝试插入、删除、替换三种操作。

\textbf{暴力过程。} 以 word1=ab、word2=ac 为例，末尾 b 和 c 不同。替换 b 成 c 来自 dp[1][1]+1；删除 b 来自 dp[1][2]+1；插入 c 来自 dp[2][1]+1。

\textbf{重复工作。} 题面模型：二维状态表示两个前缀之间的最少编辑操作。暴力版的慢点要落到本题：浪费在不同编辑序列会反复到达同一对前缀长度。只要 i 和 j 相同，后续最少编辑次数就相同。后面的优化只消掉这类重复工作，不能改变答案集合。

\textbf{优化条件。} 本题的关键观察是：末尾字符相同继承左上；不同则枚举插入、删除、替换。这句话必须能翻译成可维护状态；如果题目条件改变后观察不再成立，就不能继续套原来的写法。

\textbf{相邻状态。} 规律是每个格子枚举最后一次编辑操作，而不是凭感觉填三邻居最小。把这个特例继续拆开看：先问暴力搜索里哪些不同路径到达同一个后续问题，再给这个后续问题命名为状态。状态必须包含未来决策需要的全部信息，转移只从更小且已经算清的状态来。

\textbf{状态复用。} dp[i][j] 表示 word1 前 i 个字符变成 word2 前 j 个字符的最少操作数。按前缀长度从小到大填表。手算时只改被新输入影响的那一小部分，而不是回头重算完整候选。

\textbf{指针和字段。} 状态字段要能说清：i/j 是前缀长度；dp[i-1][j] 对应删除，dp[i][j-1] 对应插入，dp[i-1][j-1] 对应替换或字符相等继承。手算时按这个协议跟踪：空串到 ac 需要插入 2 次；ab 到空串需要删除 2 次。内部格子按末尾字符是否相等决定转移。

\textbf{代码落点。} 关键是第一行 dp[0][j]=j、第一列 dp[i][0]=i。字符相等时不要额外加一。关键代码行必须能逐句翻译成“旧状态怎样变成新状态”，否则就只是模板。

\textbf{正确性。} 任意最优编辑序列的最后一步只有插入、删除、替换或不操作四类，转移枚举了全部最后一步。

\textbf{边界实验。} 先用空串、完全相同、完全不同、操作代价不同做反例检查。实验时覆盖空串、完全相同、完全不同、操作代价不同，先打印完整 DP 表，再比较空间压缩版本。若压缩版不同，优先检查遍历顺序和旧值保存。

\textbf{复杂度变化。} 暴力递归指数级；二维 DP O(mn) 时间和空间，可滚动到 O(min(m, n)) 空间。

\textbf{迁移追问。} 如果题目改成“若只允许插入删除，问题退化到和最长公共子序列相关”，先判断原状态是否仍然覆盖新答案集合，再决定是微调边界、改 value 语义，还是换算法。

\noindent\textbf{LeetCode 97 交错字符串。}

\textbf{题目说明。} LeetCode 97 交错字符串 的题面入口要先还原成具体任务：目标串前缀由两个来源字符串的前缀交错组成。

\textbf{样例入口。} 拿 s1=a、s2=b、s3=ab 手算，最后一个字符 b 可以来自 s2；拿 s1=a、s2=a、s3=aa，两个来源都能匹配，不能贪心固定选一个。

\textbf{暴力基线。} 慢版本先写递归搜索树：节点表示处理位置、剩余资源和当前选择。只要不同路径反复到达同一个节点，就说明需要把这个节点命名成 DP 状态。放到本题就是：先按“目标串前缀由两个来源字符串的前缀交错组成”完整试慢版本；样例里已经能看到“规律是 dp[i][j] 表示两个前缀能否组成目标前 i+j 个字符，转移同时尝试来自 s1 和 s2 的最后一步”，所以优化前必须先能说清慢版本枚举了哪些候选。

\textbf{暴力过程。} 拿 s1=a、s2=b、s3=ab 手算，最后一个字符 b 可以来自 s2；拿 s1=a、s2=a、s3=aa，两个来源都能匹配，不能贪心固定选一个。

\textbf{重复工作。} 题面模型：目标串前缀由两个来源字符串的前缀交错组成。暴力版的慢点要落到本题：慢版本会在“目标串前缀由两个来源字符串的前缀交错组成”的选择树里反复到达相同参数。样例规律是“规律是 dp[i][j] 表示两个前缀能否组成目标前 i+j 个字符，转移同时尝试来自 s1 和 s2 的最后一步”，说明这个参数组合可以命名成状态，算过之后后面直接复用。后面的优化只消掉这类重复工作，不能改变答案集合。

\textbf{优化条件。} 本题的关键观察是：状态 dp[i][j] 表示 s1 前 i 个字符和 s2 前 j 个字符能否组成 s3 前 i 加 j 个字符。这句话必须能翻译成可维护状态；如果题目条件改变后观察不再成立，就不能继续套原来的写法。

\textbf{相邻状态。} 规律是 dp[i][j] 表示两个前缀能否组成目标前 i+j 个字符，转移同时尝试来自 s1 和 s2 的最后一步。把这个特例继续拆开看：先问暴力搜索里哪些不同路径到达同一个后续问题，再给这个后续问题命名为状态。状态必须包含未来决策需要的全部信息，转移只从更小且已经算清的状态来。

\textbf{状态复用。} 把重复递归节点命名为状态；遍历顺序只是在保证依赖状态已经算完。本题落点是：规律是 dp[i][j] 表示两个前缀能否组成目标前 i+j 个字符，转移同时尝试来自 s1 和 s2 的最后一步。手算时只改被新输入影响的那一小部分，而不是回头重算完整候选。

\textbf{指针和字段。} 状态字段要能说清：状态下标、状态值语义、初始化、转移来源、遍历顺序；空间压缩时还要指出读到的是旧值还是新值。手算时按这个协议跟踪：拿 s1=a、s2=b、s3=ab 手算，最后一个字符 b 可以来自 s2；拿 s1=a、s2=a、s3=aa，两个来源都能匹配，不能贪心固定选一个。然后按协议复查：手算时先写一个很小的 DP 表或记忆化搜索记录。每个格子旁边写它来自哪些更小状态。

\textbf{代码落点。} 先写未压缩版本校验转移；压缩前后用小表对拍；不可达哨兵参与加法前要检查溢出。关键代码行必须能逐句翻译成“旧状态怎样变成新状态”，否则就只是模板。

\textbf{正确性。} 证明从最后一步开始：任意最优解的最后一步都在转移枚举中，转移来源已经由更小状态正确计算。

\textbf{边界实验。} 先用总长度不等、第一行第一列初始化、两个来源字符都能匹配时不能贪心只选一个做反例检查。实验时覆盖总长度不等、第一行第一列初始化、两个来源字符都能匹配时不能贪心只选一个，先打印完整 DP 表，再比较空间压缩版本。若压缩版不同，优先检查遍历顺序和旧值保存。

\textbf{复杂度变化。} 暴力递归会重复计算相同子问题，常见指数级；DP 把每个状态算一次，时间是状态数乘单个转移成本，空间是状态表大小或压缩后的大小。

\textbf{迁移追问。} 如果题目改成“若要统计交错方案数，布尔状态要改成计数状态并处理大数或取模”，先判断原状态是否仍然覆盖新答案集合，再决定是微调边界、改 value 语义，还是换算法。

\noindent\textbf{LeetCode 115 不同的子序列。}

\textbf{题目说明。} LeetCode 115 不同的子序列 的题面入口要先还原成具体任务：从源串中删除若干字符后形成目标串，本质是两个前缀之间的计数问题。

\textbf{样例入口。} 拿 s=bab、t=ba 手算。扫描到最后一个 b 时，它不能匹配 t 的 a，只能继承不用它的方案；扫描到 a 时，有两类方案：用这个 a 匹配 t 的末尾，或不用它。

\textbf{暴力基线。} 慢版本先写递归搜索树：节点表示处理位置、剩余资源和当前选择。只要不同路径反复到达同一个节点，就说明需要把这个节点命名成 DP 状态。放到本题就是：先按“从源串中删除若干字符后形成目标串，本质是两个前缀之间的计数问题”完整试慢版本；样例里已经能看到“规律是字符相等时计数相加，不等时只继承不用源字符的方案”，所以优化前必须先能说清慢版本枚举了哪些候选。

\textbf{暴力过程。} 拿 s=bab、t=ba 手算。扫描到最后一个 b 时，它不能匹配 t 的 a，只能继承不用它的方案；扫描到 a 时，有两类方案：用这个 a 匹配 t 的末尾，或不用它。

\textbf{重复工作。} 题面模型：从源串中删除若干字符后形成目标串，本质是两个前缀之间的计数问题。暴力版的慢点要落到本题：慢版本会在“从源串中删除若干字符后形成目标串，本质是两个前缀之间的计数问题”的选择树里反复到达相同参数。样例规律是“规律是字符相等时计数相加，不等时只继承不用源字符的方案”，说明这个参数组合可以命名成状态，算过之后后面直接复用。后面的优化只消掉这类重复工作，不能改变答案集合。

\textbf{优化条件。} 本题的关键观察是：状态 dp[i][j] 表示源串前 i 个字符中有多少个子序列等于目标前 j 个字符。这句话必须能翻译成可维护状态；如果题目条件改变后观察不再成立，就不能继续套原来的写法。

\textbf{相邻状态。} 规律是字符相等时计数相加，不等时只继承不用源字符的方案。把这个特例继续拆开看：先问暴力搜索里哪些不同路径到达同一个后续问题，再给这个后续问题命名为状态。状态必须包含未来决策需要的全部信息，转移只从更小且已经算清的状态来。

\textbf{状态复用。} 把重复递归节点命名为状态；遍历顺序只是在保证依赖状态已经算完。本题落点是：规律是字符相等时计数相加，不等时只继承不用源字符的方案。手算时只改被新输入影响的那一小部分，而不是回头重算完整候选。

\textbf{指针和字段。} 状态字段要能说清：状态下标、状态值语义、初始化、转移来源、遍历顺序；空间压缩时还要指出读到的是旧值还是新值。手算时按这个协议跟踪：拿 s=bab、t=ba 手算。扫描到最后一个 b 时，它不能匹配 t 的 a，只能继承不用它的方案；扫描到 a 时，有两类方案：用这个 a 匹配 t 的末尾，或不用它。然后按协议复查：手算时先写一个很小的 DP 表或记忆化搜索记录。每个格子旁边写它来自哪些更小状态。

\textbf{代码落点。} 先写未压缩版本校验转移；压缩前后用小表对拍；不可达哨兵参与加法前要检查溢出。关键代码行必须能逐句翻译成“旧状态怎样变成新状态”，否则就只是模板。

\textbf{正确性。} 证明从最后一步开始：任意最优解的最后一步都在转移枚举中，转移来源已经由更小状态正确计算。

\textbf{边界实验。} 先用空目标串计数为一、源串为空、重复字符导致计数增长、计数可能超过 int做反例检查。实验时覆盖空目标串计数为一、源串为空、重复字符导致计数增长、计数可能超过 int，先打印完整 DP 表，再比较空间压缩版本。若压缩版不同，优先检查遍历顺序和旧值保存。

\textbf{复杂度变化。} 暴力递归会重复计算相同子问题，常见指数级；DP 把每个状态算一次，时间是状态数乘单个转移成本，空间是状态表大小或压缩后的大小。

\textbf{迁移追问。} 如果题目改成“若只问是否存在目标子序列，可以用双指针，不需要二维计数 DP”，先判断原状态是否仍然覆盖新答案集合，再决定是微调边界、改 value 语义，还是换算法。

\noindent\textbf{LeetCode 121 买卖股票的最佳时机。}

\textbf{题目说明。} LeetCode 121 买卖股票的最佳时机 的题面入口要先还原成具体任务：卖出日固定时，最佳买入日是此前最低价格。

\textbf{样例入口。} 拿价格 [7, 1, 5] 手算。站在 7 不能卖出获利；站在 1 更新最低买入价；站在 5 时利润是 5-1。

\textbf{暴力基线。} 暴力枚举买入日 i 和卖出日 j，要求 i<j，计算 prices[j]-prices[i] 的最大值。

\textbf{暴力过程。} 以 [7, 1, 5, 3, 6, 4] 为例，若卖出日在价格 6，最好的买入日是此前最低价 1，利润 5。

\textbf{重复工作。} 题面模型：卖出日固定时，最佳买入日是此前最低价格。暴力版的慢点要落到本题：浪费在每个卖出日都回头找最低买入价。扫描过程中只需要维护历史最低价。后面的优化只消掉这类重复工作，不能改变答案集合。

\textbf{优化条件。} 本题的关键观察是：扫描维护历史最低价和最大利润，一次交易无需复杂状态机。这句话必须能翻译成可维护状态；如果题目条件改变后观察不再成立，就不能继续套原来的写法。

\textbf{相邻状态。} 规律是固定卖出日后，最佳买入日只需要历史最低价，状态不是所有买卖对。把这个特例继续拆开看：先问暴力搜索里哪些不同路径到达同一个后续问题，再给这个后续问题命名为状态。状态必须包含未来决策需要的全部信息，转移只从更小且已经算清的状态来。

\textbf{状态复用。} min\_price 保存当前日前的最低价格，best 保存最大利润。每天先用当前价格尝试卖出，再更新 min\_price。手算时只改被新输入影响的那一小部分，而不是回头重算完整候选。

\textbf{指针和字段。} 状态字段要能说清：price 是当前卖出价候选，min\_price 是历史最优买入价，best 是已见最大利润。手算时按这个协议跟踪：扫描到 5 时 min\_price=1，best=4；扫描到 6 时利润 5，更新 best。

\textbf{代码落点。} 核心顺序是先用当前价格尝试卖出更新 best，再用当前价格更新 min\_price。若先更新最低价，同一天买卖利润为 0，不影响答案但语义要讲清楚。关键代码行必须能逐句翻译成“旧状态怎样变成新状态”，否则就只是模板。

\textbf{正确性。} 任意最优交易按卖出日分类；固定卖出日时，最佳买入日一定是它之前最低价格，状态完整覆盖这个选择。

\textbf{边界实验。} 先用单日价格、一直下跌、一直上涨、价格相同做反例检查。实验时覆盖单日价格、一直下跌、一直上涨、价格相同，先打印完整 DP 表，再比较空间压缩版本。若压缩版不同，优先检查遍历顺序和旧值保存。

\textbf{复杂度变化。} 暴力 O(\ensuremath{n^{2}})，一次扫描 O(n) 时间、O(1) 空间。

\textbf{迁移追问。} 如果题目改成“多次交易、冷冻期和手续费版本都要扩展成状态机”，先判断原状态是否仍然覆盖新答案集合，再决定是微调边界、改 value 语义，还是换算法。

\noindent\textbf{LeetCode 123 买卖股票 III。}

\textbf{题目说明。} LeetCode 123 买卖股票 III 的题面入口要先还原成具体任务：最多两次交易，状态需要包含交易次数和持有状态。

\textbf{样例入口。} 拿 [3, 3, 5, 0, 0, 3, 1, 4] 手算，只保存一个最大利润不够，因为第二次买入依赖第一次卖出后的余额。

\textbf{暴力基线。} 慢版本先写递归搜索树：节点表示处理位置、剩余资源和当前选择。只要不同路径反复到达同一个节点，就说明需要把这个节点命名成 DP 状态。放到本题就是：先按“最多两次交易，状态需要包含交易次数和持有状态”完整试慢版本；样例里已经能看到“规律是维护 buy1、sell1、buy2、sell2 四个状态，每天都表示到今天为止对应状态的最优值；两次交易不能重叠”，所以优化前必须先能说清慢版本枚举了哪些候选。

\textbf{暴力过程。} 拿 [3, 3, 5, 0, 0, 3, 1, 4] 手算，只保存一个最大利润不够，因为第二次买入依赖第一次卖出后的余额。

\textbf{重复工作。} 题面模型：最多两次交易，状态需要包含交易次数和持有状态。暴力版的慢点要落到本题：慢版本会在“最多两次交易，状态需要包含交易次数和持有状态”的选择树里反复到达相同参数。样例规律是“规律是维护 buy1、sell1、buy2、sell2 四个状态，每天都表示到今天为止对应状态的最优值；两次交易不能重叠”，说明这个参数组合可以命名成状态，算过之后后面直接复用。后面的优化只消掉这类重复工作，不能改变答案集合。

\textbf{优化条件。} 本题的关键观察是：维护 buy1、sell1、buy2、sell2 四个状态，按天更新。这句话必须能翻译成可维护状态；如果题目条件改变后观察不再成立，就不能继续套原来的写法。

\textbf{相邻状态。} 规律是维护 buy1、sell1、buy2、sell2 四个状态，每天都表示到今天为止对应状态的最优值；两次交易不能重叠。把这个特例继续拆开看：先问暴力搜索里哪些不同路径到达同一个后续问题，再给这个后续问题命名为状态。状态必须包含未来决策需要的全部信息，转移只从更小且已经算清的状态来。

\textbf{状态复用。} 把重复递归节点命名为状态；遍历顺序只是在保证依赖状态已经算完。本题落点是：规律是维护 buy1、sell1、buy2、sell2 四个状态，每天都表示到今天为止对应状态的最优值；两次交易不能重叠。手算时只改被新输入影响的那一小部分，而不是回头重算完整候选。

\textbf{指针和字段。} 状态字段要能说清：状态下标、状态值语义、初始化、转移来源、遍历顺序；空间压缩时还要指出读到的是旧值还是新值。手算时按这个协议跟踪：拿 [3, 3, 5, 0, 0, 3, 1, 4] 手算，只保存一个最大利润不够，因为第二次买入依赖第一次卖出后的余额。然后按协议复查：手算时先写一个很小的 DP 表或记忆化搜索记录。每个格子旁边写它来自哪些更小状态。

\textbf{代码落点。} 先写未压缩版本校验转移；压缩前后用小表对拍；不可达哨兵参与加法前要检查溢出。关键代码行必须能逐句翻译成“旧状态怎样变成新状态”，否则就只是模板。

\textbf{正确性。} 证明从最后一步开始：任意最优解的最后一步都在转移枚举中，转移来源已经由更小状态正确计算。

\textbf{边界实验。} 先用价格为空、只够一次交易、两次交易间不能重叠做反例检查。实验时覆盖价格为空、只够一次交易、两次交易间不能重叠，先打印完整 DP 表，再比较空间压缩版本。若压缩版不同，优先检查遍历顺序和旧值保存。

\textbf{复杂度变化。} 暴力递归会重复计算相同子问题，常见指数级；DP 把每个状态算一次，时间是状态数乘单个转移成本，空间是状态表大小或压缩后的大小。

\textbf{迁移追问。} 如果题目改成“最多 k 次交易是它的泛化，k 很大时退化为无限次交易”，先判断原状态是否仍然覆盖新答案集合，再决定是微调边界、改 value 语义，还是换算法。

\noindent\textbf{LeetCode 139 单词拆分。}

\textbf{题目说明。} LeetCode 139 单词拆分 的题面入口要先还原成具体任务：状态表示前缀能否由字典单词拼出。

\textbf{样例入口。} 拿 s=leetcode、字典 [leet, code] 手算，dp[4] 因为 leet 可拆而为真，dp[8] 又因为 dp[4] 且 code 在字典里为真。

\textbf{暴力基线。} 暴力递归从字符串开头切一个字典单词，剩余后缀继续递归；尝试所有切法。

\textbf{暴力过程。} 以 leetcode 和字典 [leet, code] 为例，前缀 leet 在字典中，剩余 code 也在字典中，所以可拆。

\textbf{重复工作。} 题面模型：状态表示前缀能否由字典单词拼出。暴力版的慢点要落到本题：浪费在同一个后缀会由不同切法反复判断。例如某个下标 i 之后能否拆分，只取决于 i，不取决于前面怎么切。后面的优化只消掉这类重复工作，不能改变答案集合。

\textbf{优化条件。} 本题的关键观察是：枚举最后一个单词的起点，左侧前缀可达且子串在字典中则当前可达。这句话必须能翻译成可维护状态；如果题目条件改变后观察不再成立，就不能继续套原来的写法。

\textbf{相邻状态。} 规律是 dp[i] 表示前 i 个字符能否拆分，转移枚举最后一个单词起点。把这个特例继续拆开看：先问暴力搜索里哪些不同路径到达同一个后续问题，再给这个后续问题命名为状态。状态必须包含未来决策需要的全部信息，转移只从更小且已经算清的状态来。

\textbf{状态复用。} dp[i] 表示 s[0..i) 是否能被字典拆分。若存在 j<i，使 dp[j] 为真且 s[j..i) 在字典中，则 dp[i]=true。手算时只改被新输入影响的那一小部分，而不是回头重算完整候选。

\textbf{指针和字段。} 状态字段要能说清：i 是当前前缀右端，j 是最后一个单词起点，word\_set 用于 O(1) 查询子串是否为单词。手算时按这个协议跟踪：leetcode 中 dp[4] 因 leet 为真；随后 i=8 时 j=4，dp[4] 为真且 code 在字典中，所以 dp[8]=true。

\textbf{代码落点。} 关键是 dp[0]=true 表示空前缀可拆。热循环里 substr 会复制字符串，工程版可按最大单词长度限制 j 或用 Trie。关键代码行必须能逐句翻译成“旧状态怎样变成新状态”，否则就只是模板。

\textbf{正确性。} 任意合法拆分都有最后一个单词 s[j..i)，其左侧前缀必须可拆；转移枚举了所有可能最后单词起点。

\textbf{边界实验。} 先用空串、重复单词、长单词、substr 复制成本做反例检查。实验时覆盖空串、重复单词、长单词、substr 复制成本，先打印完整 DP 表，再比较空间压缩版本。若压缩版不同，优先检查遍历顺序和旧值保存。

\textbf{复杂度变化。} 暴力指数级；DP 有 O(\ensuremath{n^{2}}) 个切分检查，若 substr 计入成本可到 O(\ensuremath{n^{3}})，空间 O(n+字典)。

\textbf{迁移追问。} 如果题目改成“若要输出所有句子，需要 DFS 加记忆化或前驱列表”，先判断原状态是否仍然覆盖新答案集合，再决定是微调边界、改 value 语义，还是换算法。

\noindent\textbf{LeetCode 152 乘积最大子数组。}

\textbf{题目说明。} LeetCode 152 乘积最大子数组 的题面入口要先还原成具体任务：负数会让最大乘积和最小乘积互换。

\textbf{样例入口。} 拿 [-2, 3, -4] 手算。到 3 时最大乘积是 3、最小是 -6；再乘 -4 后，原来的最小 -6 反而变成最大 24。

\textbf{暴力基线。} 慢版本先写递归搜索树：节点表示处理位置、剩余资源和当前选择。只要不同路径反复到达同一个节点，就说明需要把这个节点命名成 DP 状态。放到本题就是：先按“负数会让最大乘积和最小乘积互换”完整试慢版本；样例里已经能看到“规律是负数会交换最大和最小，状态必须同时保存以当前位置结尾的最大乘积和最小乘积”，所以优化前必须先能说清慢版本枚举了哪些候选。

\textbf{暴力过程。} 拿 [-2, 3, -4] 手算。到 3 时最大乘积是 3、最小是 -6；再乘 -4 后，原来的最小 -6 反而变成最大 24。

\textbf{重复工作。} 题面模型：负数会让最大乘积和最小乘积互换。暴力版的慢点要落到本题：慢版本会在“负数会让最大乘积和最小乘积互换”的选择树里反复到达相同参数。样例规律是“规律是负数会交换最大和最小，状态必须同时保存以当前位置结尾的最大乘积和最小乘积”，说明这个参数组合可以命名成状态，算过之后后面直接复用。后面的优化只消掉这类重复工作，不能改变答案集合。

\textbf{优化条件。} 本题的关键观察是：同时维护以当前位置结尾的最大和最小乘积。这句话必须能翻译成可维护状态；如果题目条件改变后观察不再成立，就不能继续套原来的写法。

\textbf{相邻状态。} 规律是负数会交换最大和最小，状态必须同时保存以当前位置结尾的最大乘积和最小乘积。把这个特例继续拆开看：先问暴力搜索里哪些不同路径到达同一个后续问题，再给这个后续问题命名为状态。状态必须包含未来决策需要的全部信息，转移只从更小且已经算清的状态来。

\textbf{状态复用。} 把重复递归节点命名为状态；遍历顺序只是在保证依赖状态已经算完。本题落点是：规律是负数会交换最大和最小，状态必须同时保存以当前位置结尾的最大乘积和最小乘积。手算时只改被新输入影响的那一小部分，而不是回头重算完整候选。

\textbf{指针和字段。} 状态字段要能说清：状态下标、状态值语义、初始化、转移来源、遍历顺序；空间压缩时还要指出读到的是旧值还是新值。手算时按这个协议跟踪：拿 [-2, 3, -4] 手算。到 3 时最大乘积是 3、最小是 -6；再乘 -4 后，原来的最小 -6 反而变成最大 24。然后按协议复查：手算时先写一个很小的 DP 表或记忆化搜索记录。每个格子旁边写它来自哪些更小状态。

\textbf{代码落点。} 先写未压缩版本校验转移；压缩前后用小表对拍；不可达哨兵参与加法前要检查溢出。关键代码行必须能逐句翻译成“旧状态怎样变成新状态”，否则就只是模板。

\textbf{正确性。} 证明从最后一步开始：任意最优解的最后一步都在转移枚举中，转移来源已经由更小状态正确计算。

\textbf{边界实验。} 先用零、负数个数奇偶、单元素、乘积溢出做反例检查。实验时覆盖零、负数个数奇偶、单元素、乘积溢出，先打印完整 DP 表，再比较空间压缩版本。若压缩版不同，优先检查遍历顺序和旧值保存。

\textbf{复杂度变化。} 暴力递归会重复计算相同子问题，常见指数级；DP 把每个状态算一次，时间是状态数乘单个转移成本，空间是状态表大小或压缩后的大小。

\textbf{迁移追问。} 如果题目改成“如果要求返回区间，同步记录状态来源”，先判断原状态是否仍然覆盖新答案集合，再决定是微调边界、改 value 语义，还是换算法。

\noindent\textbf{LeetCode 188 买卖股票 IV。}

\textbf{题目说明。} LeetCode 188 买卖股票 IV 的题面入口要先还原成具体任务：最多 k 次交易让状态必须同时包含交易次数和是否持有股票。

\textbf{样例入口。} 拿 k=2、价格 [3, 2, 6] 手算。某天结束时只说最大利润不够，因为未来能否卖出取决于是否持股，也取决于已经完成几次交易。

\textbf{暴力基线。} 慢版本先写递归搜索树：节点表示处理位置、剩余资源和当前选择。只要不同路径反复到达同一个节点，就说明需要把这个节点命名成 DP 状态。放到本题就是：先按“最多 k 次交易让状态必须同时包含交易次数和是否持有股票”完整试慢版本；样例里已经能看到“规律是状态必须包含交易次数和持有状态；k 很大时才退化为无限次交易”，所以优化前必须先能说清慢版本枚举了哪些候选。

\textbf{暴力过程。} 拿 k=2、价格 [3, 2, 6] 手算。某天结束时只说最大利润不够，因为未来能否卖出取决于是否持股，也取决于已经完成几次交易。

\textbf{重复工作。} 题面模型：最多 k 次交易让状态必须同时包含交易次数和是否持有股票。暴力版的慢点要落到本题：慢版本会在“最多 k 次交易让状态必须同时包含交易次数和是否持有股票”的选择树里反复到达相同参数。样例规律是“规律是状态必须包含交易次数和持有状态；k 很大时才退化为无限次交易”，说明这个参数组合可以命名成状态，算过之后后面直接复用。后面的优化只消掉这类重复工作，不能改变答案集合。

\textbf{优化条件。} 本题的关键观察是：维护每个交易次数下的 buy 和 sell 状态；当 k 足够大时可退化为无限次交易。这句话必须能翻译成可维护状态；如果题目条件改变后观察不再成立，就不能继续套原来的写法。

\textbf{相邻状态。} 规律是状态必须包含交易次数和持有状态；k 很大时才退化为无限次交易。把这个特例继续拆开看：先问暴力搜索里哪些不同路径到达同一个后续问题，再给这个后续问题命名为状态。状态必须包含未来决策需要的全部信息，转移只从更小且已经算清的状态来。

\textbf{状态复用。} 把重复递归节点命名为状态；遍历顺序只是在保证依赖状态已经算完。本题落点是：规律是状态必须包含交易次数和持有状态；k 很大时才退化为无限次交易。手算时只改被新输入影响的那一小部分，而不是回头重算完整候选。

\textbf{指针和字段。} 状态字段要能说清：状态下标、状态值语义、初始化、转移来源、遍历顺序；空间压缩时还要指出读到的是旧值还是新值。手算时按这个协议跟踪：拿 k=2、价格 [3, 2, 6] 手算。某天结束时只说最大利润不够，因为未来能否卖出取决于是否持股，也取决于已经完成几次交易。然后按协议复查：手算时先写一个很小的 DP 表或记忆化搜索记录。每个格子旁边写它来自哪些更小状态。

\textbf{代码落点。} 先写未压缩版本校验转移；压缩前后用小表对拍；不可达哨兵参与加法前要检查溢出。关键代码行必须能逐句翻译成“旧状态怎样变成新状态”，否则就只是模板。

\textbf{正确性。} 证明从最后一步开始：任意最优解的最后一步都在转移枚举中，转移来源已经由更小状态正确计算。

\textbf{边界实验。} 先用k 为零、价格天数很少、k 大于等于 n/2、更新顺序造成同一天状态污染做反例检查。实验时覆盖k 为零、价格天数很少、k 大于等于 n/2、更新顺序造成同一天状态污染，先打印完整 DP 表，再比较空间压缩版本。若压缩版不同，优先检查遍历顺序和旧值保存。

\textbf{复杂度变化。} 暴力递归会重复计算相同子问题，常见指数级；DP 把每个状态算一次，时间是状态数乘单个转移成本，空间是状态表大小或压缩后的大小。

\textbf{迁移追问。} 如果题目改成“手续费、冷冻期和最多两次交易都是同一状态机框架的不同约束”，先判断原状态是否仍然覆盖新答案集合，再决定是微调边界、改 value 语义，还是换算法。

\noindent\textbf{LeetCode 198 打家劫舍。}

\textbf{题目说明。} LeetCode 198 打家劫舍 的题面入口要先还原成具体任务：每间房只有偷和不偷两种结局，偷当前就不能偷前一间。

\textbf{样例入口。} 拿 [2, 7, 9] 手算，到房屋 9 时只有两类合法结局：不偷它，答案沿用前一间；偷它，就只能接前两间的最优。

\textbf{暴力基线。} 暴力递归每间房偷或不偷，偷了当前房就不能偷下一间，最后取最大金额。

\textbf{暴力过程。} 以 nums=[2, 7, 9] 为例，到房屋 9 时，合法结局只有两类：不偷 9 沿用前一间最优，偷 9 就只能接前两间最优。

\textbf{重复工作。} 题面模型：每间房只有偷和不偷两种结局，偷当前就不能偷前一间。暴力版的慢点要落到本题：浪费在递归中同一个下标之后的最优收益被反复计算。未来只取决于当前处理到哪一间。后面的优化只消掉这类重复工作，不能改变答案集合。

\textbf{优化条件。} 本题的关键观察是：滚动保存前一状态和前两状态。这句话必须能翻译成可维护状态；如果题目条件改变后观察不再成立，就不能继续套原来的写法。

\textbf{相邻状态。} 规律是 dp[i]=max(dp[i-1], dp[i-2]+nums[i])，相邻限制来自题面而不是公式本身。把这个特例继续拆开看：先问暴力搜索里哪些不同路径到达同一个后续问题，再给这个后续问题命名为状态。状态必须包含未来决策需要的全部信息，转移只从更小且已经算清的状态来。

\textbf{状态复用。} dp[i] 表示考虑前 i 间房的最大收益。转移为 max(dp[i-1], dp[i-2]+nums[i-1])，可压成两个变量。手算时只改被新输入影响的那一小部分，而不是回头重算完整候选。

\textbf{指针和字段。} 状态字段要能说清：previous\_one 是 dp[i-1]，previous\_two 是 dp[i-2]，money 是当前房屋金额。手算时按这个协议跟踪：[2, 7, 9] 中前两步最优到 7；处理 9 时比较不偷得 7、偷得 2+9=11，更新为 11。

\textbf{代码落点。} 关键是滚动更新顺序：先算 current=max(previous\_one, previous\_two+money)，再 previous\_two=previous\_one，previous\_one=current。关键代码行必须能逐句翻译成“旧状态怎样变成新状态”，否则就只是模板。

\textbf{正确性。} 最优方案对第 i 间房只有偷或不偷两类，二者互斥且覆盖全部合法方案。

\textbf{边界实验。} 先用空数组、单间、全零、金额很大做反例检查。实验时覆盖空数组、单间、全零、金额很大，先打印完整 DP 表，再比较空间压缩版本。若压缩版不同，优先检查遍历顺序和旧值保存。

\textbf{复杂度变化。} 暴力递归指数级；DP O(n) 时间，滚动 O(1) 空间。

\textbf{迁移追问。} 如果题目改成“环形房屋要拆成不含首和不含尾两个线性问题”，先判断原状态是否仍然覆盖新答案集合，再决定是微调边界、改 value 语义，还是换算法。

\noindent\textbf{LeetCode 221 最大正方形。}

\textbf{题目说明。} LeetCode 221 最大正方形 的题面入口要先还原成具体任务：以当前位置为右下角的最大正方形由上、左、左上三个状态限制。

\textbf{样例入口。} 拿 2x2 全 1 方块手算，右下角能形成边长 2，因为上、左、左上三个格子都至少能形成边长 1。若其中任一方向为 0，只能形成边长 1。

\textbf{暴力基线。} 慢版本先写递归搜索树：节点表示处理位置、剩余资源和当前选择。只要不同路径反复到达同一个节点，就说明需要把这个节点命名成 DP 状态。放到本题就是：先按“以当前位置为右下角的最大正方形由上、左、左上三个状态限制”完整试慢版本；样例里已经能看到“规律是 dp[row][col]=min(上, 左, 左上)+1，状态值是以当前格为右下角的最大正方形边长”，所以优化前必须先能说清慢版本枚举了哪些候选。

\textbf{暴力过程。} 拿 2x2 全 1 方块手算，右下角能形成边长 2，因为上、左、左上三个格子都至少能形成边长 1。若其中任一方向为 0，只能形成边长 1。

\textbf{重复工作。} 题面模型：以当前位置为右下角的最大正方形由上、左、左上三个状态限制。暴力版的慢点要落到本题：慢版本会在“以当前位置为右下角的最大正方形由上、左、左上三个状态限制”的选择树里反复到达相同参数。样例规律是“规律是 dp[row][col]=min(上, 左, 左上)+1，状态值是以当前格为右下角的最大正方形边长”，说明这个参数组合可以命名成状态，算过之后后面直接复用。后面的优化只消掉这类重复工作，不能改变答案集合。

\textbf{优化条件。} 本题的关键观察是：当前位置为一时取三者最小加一，否则为零。这句话必须能翻译成可维护状态；如果题目条件改变后观察不再成立，就不能继续套原来的写法。

\textbf{相邻状态。} 规律是 dp[row][col]=min(上, 左, 左上)+1，状态值是以当前格为右下角的最大正方形边长。把这个特例继续拆开看：先问暴力搜索里哪些不同路径到达同一个后续问题，再给这个后续问题命名为状态。状态必须包含未来决策需要的全部信息，转移只从更小且已经算清的状态来。

\textbf{状态复用。} 把重复递归节点命名为状态；遍历顺序只是在保证依赖状态已经算完。本题落点是：规律是 dp[row][col]=min(上, 左, 左上)+1，状态值是以当前格为右下角的最大正方形边长。手算时只改被新输入影响的那一小部分，而不是回头重算完整候选。

\textbf{指针和字段。} 状态字段要能说清：状态下标、状态值语义、初始化、转移来源、遍历顺序；空间压缩时还要指出读到的是旧值还是新值。手算时按这个协议跟踪：拿 2x2 全 1 方块手算，右下角能形成边长 2，因为上、左、左上三个格子都至少能形成边长 1。若其中任一方向为 0，只能形成边长 1。然后按协议复查：手算时先写一个很小的 DP 表或记忆化搜索记录。每个格子旁边写它来自哪些更小状态。

\textbf{代码落点。} 先写未压缩版本校验转移；压缩前后用小表对拍；不可达哨兵参与加法前要检查溢出。关键代码行必须能逐句翻译成“旧状态怎样变成新状态”，否则就只是模板。

\textbf{正确性。} 证明从最后一步开始：任意最优解的最后一步都在转移枚举中，转移来源已经由更小状态正确计算。

\textbf{边界实验。} 先用全零、全一、单行单列、字符一和数字一不要混淆做反例检查。实验时覆盖全零、全一、单行单列、字符一和数字一不要混淆，先打印完整 DP 表，再比较空间压缩版本。若压缩版不同，优先检查遍历顺序和旧值保存。

\textbf{复杂度变化。} 暴力递归会重复计算相同子问题，常见指数级；DP 把每个状态算一次，时间是状态数乘单个转移成本，空间是状态表大小或压缩后的大小。

\textbf{迁移追问。} 如果题目改成“最大矩形则需要按行构造柱状图再用单调栈”，先判断原状态是否仍然覆盖新答案集合，再决定是微调边界、改 value 语义，还是换算法。

\noindent\textbf{LeetCode 300 最长递增子序列。}

\textbf{题目说明。} LeetCode 300 最长递增子序列 的题面入口要先还原成具体任务：二次 DP 固定结尾，二分优化维护每个长度的最小结尾。

\textbf{样例入口。} 拿 [10, 9, 2, 5, 3] 手算。二次 DP 固定每个位置为结尾，5 可接在 2 后；优化时 tails[长度] 保存该长度递增子序列的最小尾值，尾值越小未来越容易接。

\textbf{暴力基线。} 暴力递归或 DP 枚举每个元素作为子序列结尾，向前寻找所有更小元素。

\textbf{暴力过程。} 以 [10, 9, 2, 5, 3, 7, 101, 18] 为例，二次 DP 可以算出以 7 结尾的长度来自 2、5 或 3 后面接 7。

\textbf{重复工作。} 题面模型：二次 DP 固定结尾，二分优化维护每个长度的最小结尾。暴力版的慢点要落到本题：二次 DP 慢在每个位置都回看所有历史元素。更快写法不需要保存真实序列，只保存每个长度的最小可能结尾。后面的优化只消掉这类重复工作，不能改变答案集合。

\textbf{优化条件。} 本题的关键观察是：tails 不是实际答案序列，而是未来扩展潜力表。这句话必须能翻译成可维护状态；如果题目条件改变后观察不再成立，就不能继续套原来的写法。

\textbf{相邻状态。} 规律是二分替换 tails，不代表真实序列完整内容，而是维护未来潜力。把这个特例继续拆开看：先问暴力搜索里哪些不同路径到达同一个后续问题，再给这个后续问题命名为状态。状态必须包含未来决策需要的全部信息，转移只从更小且已经算清的状态来。

\textbf{状态复用。} tails[len] 表示长度为 len+1 的递增子序列中最小结尾。对每个 x，用 lower\_bound 找第一个 >=x 的位置并替换。手算时只改被新输入影响的那一小部分，而不是回头重算完整候选。

\textbf{指针和字段。} 状态字段要能说清：tails 不是最终答案序列，而是未来扩展潜力表；tails.size() 是当前最长长度。手算时按这个协议跟踪：处理 2、5、3 时，tails 从 [2, 5] 变成 [2, 3]，长度没变，但结尾更小，未来更容易接 7。

\textbf{代码落点。} 严格递增用 lower\_bound；若题目改成非递减，要改用 upper\_bound。要还原路径需要额外 predecessor 数组。关键代码行必须能逐句翻译成“旧状态怎样变成新状态”，否则就只是模板。

\textbf{正确性。} 同样长度下，结尾越小，未来能接上的数字集合只会更多；替换 tails 不改变已知长度，只提升扩展潜力。

\textbf{边界实验。} 先用重复元素、严格递增和非递减区别、空数组做反例检查。实验时覆盖重复元素、严格递增和非递减区别、空数组，先打印完整 DP 表，再比较空间压缩版本。若压缩版不同，优先检查遍历顺序和旧值保存。

\textbf{复杂度变化。} 二次 DP O(\ensuremath{n^{2}})，tails 加二分 O(n log n) 时间、O(n) 空间。

\textbf{迁移追问。} 如果题目改成“若要还原路径，需要额外前驱数组，不能只看 tails”，先判断原状态是否仍然覆盖新答案集合，再决定是微调边界、改 value 语义，还是换算法。

\noindent\textbf{LeetCode 309 最佳买卖股票含冷冻期。}

\textbf{题目说明。} LeetCode 309 最佳买卖股票含冷冻期 的题面入口要先还原成具体任务：冷冻期让买入来源受昨天卖出状态限制。

\textbf{样例入口。} 拿 [1, 2, 3, 0, 2] 手算，第 2 天卖出后第 3 天不能立刻买入，因为有冷冻期；状态必须区分持股、刚卖出、空仓可买。

\textbf{暴力基线。} 慢版本先写递归搜索树：节点表示处理位置、剩余资源和当前选择。只要不同路径反复到达同一个节点，就说明需要把这个节点命名成 DP 状态。放到本题就是：先按“冷冻期让买入来源受昨天卖出状态限制”完整试慢版本；样例里已经能看到“规律是冷冻期把普通买卖股票拆成多个日末状态，转移不能跨过限制”，所以优化前必须先能说清慢版本枚举了哪些候选。

\textbf{暴力过程。} 拿 [1, 2, 3, 0, 2] 手算，第 2 天卖出后第 3 天不能立刻买入，因为有冷冻期；状态必须区分持股、刚卖出、空仓可买。

\textbf{重复工作。} 题面模型：冷冻期让买入来源受昨天卖出状态限制。暴力版的慢点要落到本题：慢版本会在“冷冻期让买入来源受昨天卖出状态限制”的选择树里反复到达相同参数。样例规律是“规律是冷冻期把普通买卖股票拆成多个日末状态，转移不能跨过限制”，说明这个参数组合可以命名成状态，算过之后后面直接复用。后面的优化只消掉这类重复工作，不能改变答案集合。

\textbf{优化条件。} 本题的关键观察是：状态机维护持有、刚卖出、休息三种日终状态。这句话必须能翻译成可维护状态；如果题目条件改变后观察不再成立，就不能继续套原来的写法。

\textbf{相邻状态。} 规律是冷冻期把普通买卖股票拆成多个日末状态，转移不能跨过限制。把这个特例继续拆开看：先问暴力搜索里哪些不同路径到达同一个后续问题，再给这个后续问题命名为状态。状态必须包含未来决策需要的全部信息，转移只从更小且已经算清的状态来。

\textbf{状态复用。} 把重复递归节点命名为状态；遍历顺序只是在保证依赖状态已经算完。本题落点是：规律是冷冻期把普通买卖股票拆成多个日末状态，转移不能跨过限制。手算时只改被新输入影响的那一小部分，而不是回头重算完整候选。

\textbf{指针和字段。} 状态字段要能说清：状态下标、状态值语义、初始化、转移来源、遍历顺序；空间压缩时还要指出读到的是旧值还是新值。手算时按这个协议跟踪：拿 [1, 2, 3, 0, 2] 手算，第 2 天卖出后第 3 天不能立刻买入，因为有冷冻期；状态必须区分持股、刚卖出、空仓可买。然后按协议复查：手算时先写一个很小的 DP 表或记忆化搜索记录。每个格子旁边写它来自哪些更小状态。

\textbf{代码落点。} 先写未压缩版本校验转移；压缩前后用小表对拍；不可达哨兵参与加法前要检查溢出。关键代码行必须能逐句翻译成“旧状态怎样变成新状态”，否则就只是模板。

\textbf{正确性。} 证明从最后一步开始：任意最优解的最后一步都在转移枚举中，转移来源已经由更小状态正确计算。

\textbf{边界实验。} 先用空价格、一天、连续上涨、卖出后不能次日买入做反例检查。实验时覆盖空价格、一天、连续上涨、卖出后不能次日买入，先打印完整 DP 表，再比较空间压缩版本。若压缩版不同，优先检查遍历顺序和旧值保存。

\textbf{复杂度变化。} 暴力递归会重复计算相同子问题，常见指数级；DP 把每个状态算一次，时间是状态数乘单个转移成本，空间是状态表大小或压缩后的大小。

\textbf{迁移追问。} 如果题目改成“手续费版本可以在买入或卖出转移中扣费”，先判断原状态是否仍然覆盖新答案集合，再决定是微调边界、改 value 语义，还是换算法。

\noindent\textbf{LeetCode 312 戳气球。}

\textbf{题目说明。} LeetCode 312 戳气球 的题面入口要先还原成具体任务：先戳哪个会让邻居持续变化，改成最后戳某个气球后区间边界才稳定。

\textbf{样例入口。} 拿 [3, 1, 5] 手算，若先戳 1，左右邻居会变；若改问某个区间里最后戳谁，最后一刻左右边界固定，收益可拆成左右两个子区间。

\textbf{暴力基线。} 慢版本先写递归搜索树：节点表示处理位置、剩余资源和当前选择。只要不同路径反复到达同一个节点，就说明需要把这个节点命名成 DP 状态。放到本题就是：先按“先戳哪个会让邻居持续变化，改成最后戳某个气球后区间边界才稳定”完整试慢版本；样例里已经能看到“规律是区间 DP 枚举最后戳的 mid，而不是第一个戳的气球”，所以优化前必须先能说清慢版本枚举了哪些候选。

\textbf{暴力过程。} 拿 [3, 1, 5] 手算，若先戳 1，左右邻居会变；若改问某个区间里最后戳谁，最后一刻左右边界固定，收益可拆成左右两个子区间。

\textbf{重复工作。} 题面模型：先戳哪个会让邻居持续变化，改成最后戳某个气球后区间边界才稳定。暴力版的慢点要落到本题：慢版本会在“先戳哪个会让邻居持续变化，改成最后戳某个气球后区间边界才稳定”的选择树里反复到达相同参数。样例规律是“规律是区间 DP 枚举最后戳的 mid，而不是第一个戳的气球”，说明这个参数组合可以命名成状态，算过之后后面直接复用。后面的优化只消掉这类重复工作，不能改变答案集合。

\textbf{优化条件。} 本题的关键观察是：区间 DP 定义开区间内全部戳完的最大收益，枚举最后一个被戳的气球作为切分点。这句话必须能翻译成可维护状态；如果题目条件改变后观察不再成立，就不能继续套原来的写法。

\textbf{相邻状态。} 规律是区间 DP 枚举最后戳的 mid，而不是第一个戳的气球。把这个特例继续拆开看：先问暴力搜索里哪些不同路径到达同一个后续问题，再给这个后续问题命名为状态。状态必须包含未来决策需要的全部信息，转移只从更小且已经算清的状态来。

\textbf{状态复用。} 把重复递归节点命名为状态；遍历顺序只是在保证依赖状态已经算完。本题落点是：规律是区间 DP 枚举最后戳的 mid，而不是第一个戳的气球。手算时只改被新输入影响的那一小部分，而不是回头重算完整候选。

\textbf{指针和字段。} 状态字段要能说清：状态下标、状态值语义、初始化、转移来源、遍历顺序；空间压缩时还要指出读到的是旧值还是新值。手算时按这个协议跟踪：拿 [3, 1, 5] 手算，若先戳 1，左右邻居会变；若改问某个区间里最后戳谁，最后一刻左右边界固定，收益可拆成左右两个子区间。然后按协议复查：手算时先写一个很小的 DP 表或记忆化搜索记录。每个格子旁边写它来自哪些更小状态。

\textbf{代码落点。} 先写未压缩版本校验转移；压缩前后用小表对拍；不可达哨兵参与加法前要检查溢出。关键代码行必须能逐句翻译成“旧状态怎样变成新状态”，否则就只是模板。

\textbf{正确性。} 证明从最后一步开始：任意最优解的最后一步都在转移枚举中，转移来源已经由更小状态正确计算。

\textbf{边界实验。} 先用补一的虚拟边界、区间长度遍历顺序、空区间收益为零、乘积可能较大做反例检查。实验时覆盖补一的虚拟边界、区间长度遍历顺序、空区间收益为零、乘积可能较大，先打印完整 DP 表，再比较空间压缩版本。若压缩版不同，优先检查遍历顺序和旧值保存。

\textbf{复杂度变化。} 暴力递归会重复计算相同子问题，常见指数级；DP 把每个状态算一次，时间是状态数乘单个转移成本，空间是状态表大小或压缩后的大小。

\textbf{迁移追问。} 如果题目改成“很多区间 DP 都要换成最后一步思考，避免中间操作改变边界”，先判断原状态是否仍然覆盖新答案集合，再决定是微调边界、改 value 语义，还是换算法。

\noindent\textbf{LeetCode 329 矩阵中的最长递增路径。}

\textbf{题目说明。} LeetCode 329 矩阵中的最长递增路径 的题面入口要先还原成具体任务：每个格子的答案是从它出发能走出的最长严格递增路径。

\textbf{样例入口。} 拿矩阵中 1->2->3 的小路径手算，从 1 出发会用到从 2 出发的答案；如果另一个格子也走到 2，后缀路径被重复计算。严格递增保证不会成环。

\textbf{暴力基线。} 慢版本先写递归搜索树：节点表示处理位置、剩余资源和当前选择。只要不同路径反复到达同一个节点，就说明需要把这个节点命名成 DP 状态。放到本题就是：先按“每个格子的答案是从它出发能走出的最长严格递增路径”完整试慢版本；样例里已经能看到“规律是每个格子记忆化保存从这里出发的最长路径”，所以优化前必须先能说清慢版本枚举了哪些候选。

\textbf{暴力过程。} 拿矩阵中 1->2->3 的小路径手算，从 1 出发会用到从 2 出发的答案；如果另一个格子也走到 2，后缀路径被重复计算。严格递增保证不会成环。

\textbf{重复工作。} 题面模型：每个格子的答案是从它出发能走出的最长严格递增路径。暴力版的慢点要落到本题：慢版本会在“每个格子的答案是从它出发能走出的最长严格递增路径”的选择树里反复到达相同参数。样例规律是“规律是每个格子记忆化保存从这里出发的最长路径”，说明这个参数组合可以命名成状态，算过之后后面直接复用。后面的优化只消掉这类重复工作，不能改变答案集合。

\textbf{优化条件。} 本题的关键观察是：把严格递增关系看成无环状态图，用记忆化 DFS 或按值排序的 DAG DP 复用后续格子结果。这句话必须能翻译成可维护状态；如果题目条件改变后观察不再成立，就不能继续套原来的写法。

\textbf{相邻状态。} 规律是每个格子记忆化保存从这里出发的最长路径。把这个特例继续拆开看：先问暴力搜索里哪些不同路径到达同一个后续问题，再给这个后续问题命名为状态。状态必须包含未来决策需要的全部信息，转移只从更小且已经算清的状态来。

\textbf{状态复用。} 把重复递归节点命名为状态；遍历顺序只是在保证依赖状态已经算完。本题落点是：规律是每个格子记忆化保存从这里出发的最长路径。手算时只改被新输入影响的那一小部分，而不是回头重算完整候选。

\textbf{指针和字段。} 状态字段要能说清：状态下标、状态值语义、初始化、转移来源、遍历顺序；空间压缩时还要指出读到的是旧值还是新值。手算时按这个协议跟踪：拿矩阵中 1->2->3 的小路径手算，从 1 出发会用到从 2 出发的答案；如果另一个格子也走到 2，后缀路径被重复计算。严格递增保证不会成环。然后按协议复查：手算时先写一个很小的 DP 表或记忆化搜索记录。每个格子旁边写它来自哪些更小状态。

\textbf{代码落点。} 先写未压缩版本校验转移；压缩前后用小表对拍；不可达哨兵参与加法前要检查溢出。关键代码行必须能逐句翻译成“旧状态怎样变成新状态”，否则就只是模板。

\textbf{正确性。} 证明从最后一步开始：任意最优解的最后一步都在转移枚举中，转移来源已经由更小状态正确计算。

\textbf{边界实验。} 先用平台相等不能走、四方向越界、递归深度、允许不下降时会产生环做反例检查。实验时覆盖平台相等不能走、四方向越界、递归深度、允许不下降时会产生环，先打印完整 DP 表，再比较空间压缩版本。若压缩版不同，优先检查遍历顺序和旧值保存。

\textbf{复杂度变化。} 暴力递归会重复计算相同子问题，常见指数级；DP 把每个状态算一次，时间是状态数乘单个转移成本，空间是状态表大小或压缩后的大小。

\textbf{迁移追问。} 如果题目改成“若改成求最短路径，严格递增的 DP 模型不再直接适用”，先判断原状态是否仍然覆盖新答案集合，再决定是微调边界、改 value 语义，还是换算法。

\noindent\textbf{LeetCode 486 预测赢家。}

\textbf{题目说明。} LeetCode 486 预测赢家 的题面入口要先还原成具体任务：两人从区间两端取数，当前选择会影响对手面对的剩余区间。

\textbf{样例入口。} 拿 nums=[1, 5, 2] 手算，先手若拿 1，后手会拿 5；若先手拿 2，后手仍拿 5，先手最终都会输。dp[left][right] 保存当前玩家比对手最多多拿多少分。

\textbf{暴力基线。} 慢版本先写递归搜索树：节点表示处理位置、剩余资源和当前选择。只要不同路径反复到达同一个节点，就说明需要把这个节点命名成 DP 状态。放到本题就是：先按“两人从区间两端取数，当前选择会影响对手面对的剩余区间”完整试慢版本；样例里已经能看到“规律是选左得到 nums[left]-dp[left+1][right]，选右得到 nums[right]-dp[left][right-1]”，所以优化前必须先能说清慢版本枚举了哪些候选。

\textbf{暴力过程。} 拿 nums=[1, 5, 2] 手算，先手若拿 1，后手会拿 5；若先手拿 2，后手仍拿 5，先手最终都会输。dp[left][right] 保存当前玩家比对手最多多拿多少分。

\textbf{重复工作。} 题面模型：两人从区间两端取数，当前选择会影响对手面对的剩余区间。暴力版的慢点要落到本题：慢版本会在“两人从区间两端取数，当前选择会影响对手面对的剩余区间”的选择树里反复到达相同参数。样例规律是“规律是选左得到 nums[left]-dp[left+1][right]，选右得到 nums[right]-dp[left][right-1]”，说明这个参数组合可以命名成状态，算过之后后面直接复用。后面的优化只消掉这类重复工作，不能改变答案集合。

\textbf{优化条件。} 本题的关键观察是：dp[left][right] 表示当前玩家相对对手最多能赢多少分，取左或取右后减去子区间优势。这句话必须能翻译成可维护状态；如果题目条件改变后观察不再成立，就不能继续套原来的写法。

\textbf{相邻状态。} 规律是选左得到 nums[left]-dp[left+1][right]，选右得到 nums[right]-dp[left][right-1]。把这个特例继续拆开看：先问暴力搜索里哪些不同路径到达同一个后续问题，再给这个后续问题命名为状态。状态必须包含未来决策需要的全部信息，转移只从更小且已经算清的状态来。

\textbf{状态复用。} 把重复递归节点命名为状态；遍历顺序只是在保证依赖状态已经算完。本题落点是：规律是选左得到 nums[left]-dp[left+1][right]，选右得到 nums[right]-dp[left][right-1]。手算时只改被新输入影响的那一小部分，而不是回头重算完整候选。

\textbf{指针和字段。} 状态字段要能说清：状态下标、状态值语义、初始化、转移来源、遍历顺序；空间压缩时还要指出读到的是旧值还是新值。手算时按这个协议跟踪：拿 nums=[1, 5, 2] 手算，先手若拿 1，后手会拿 5；若先手拿 2，后手仍拿 5，先手最终都会输。dp[left][right] 保存当前玩家比对手最多多拿多少分。然后按协议复查：手算时先写一个很小的 DP 表或记忆化搜索记录。每个格子旁边写它来自哪些更小状态。

\textbf{代码落点。} 先写未压缩版本校验转移；压缩前后用小表对拍；不可达哨兵参与加法前要检查溢出。关键代码行必须能逐句翻译成“旧状态怎样变成新状态”，否则就只是模板。

\textbf{正确性。} 证明从最后一步开始：任意最优解的最后一步都在转移枚举中，转移来源已经由更小状态正确计算。

\textbf{边界实验。} 先用单元素、两个元素、总分相等、负数若存在、区间遍历顺序做反例检查。实验时覆盖单元素、两个元素、总分相等、负数若存在、区间遍历顺序，先打印完整 DP 表，再比较空间压缩版本。若压缩版不同，优先检查遍历顺序和旧值保存。

\textbf{复杂度变化。} 暴力递归会重复计算相同子问题，常见指数级；DP 把每个状态算一次，时间是状态数乘单个转移成本，空间是状态表大小或压缩后的大小。

\textbf{迁移追问。} 如果题目改成“若要输出具体策略，需要记录每个区间选择的端点”，先判断原状态是否仍然覆盖新答案集合，再决定是微调边界、改 value 语义，还是换算法。

\noindent\textbf{LeetCode 509 斐波那契数。}

\textbf{题目说明。} LeetCode 509 斐波那契数 的题面入口要先还原成具体任务：递归定义会让 fib(n-1) 和 fib(n-2) 下面的状态反复计算。

\textbf{样例入口。} 拿 n=5 手算，fib(5) 会递归调用 fib(4) 和 fib(3)，而 fib(4) 内部又会调用 fib(3)，同一个状态被重复计算。

\textbf{暴力基线。} 慢版本先写递归搜索树：节点表示处理位置、剩余资源和当前选择。只要不同路径反复到达同一个节点，就说明需要把这个节点命名成 DP 状态。放到本题就是：先按“递归定义会让 fib(n-1) 和 fib(n-2) 下面的状态反复计算”完整试慢版本；样例里已经能看到“规律是从 0 和 1 开始向上滚动保存前两项，每个 fib(i) 只算一次”，所以优化前必须先能说清慢版本枚举了哪些候选。

\textbf{暴力过程。} 拿 n=5 手算，fib(5) 会递归调用 fib(4) 和 fib(3)，而 fib(4) 内部又会调用 fib(3)，同一个状态被重复计算。

\textbf{重复工作。} 题面模型：递归定义会让 fib(n-1) 和 fib(n-2) 下面的状态反复计算。暴力版的慢点要落到本题：慢版本会在“递归定义会让 fib(n-1) 和 fib(n-2) 下面的状态反复计算”的选择树里反复到达相同参数。样例规律是“规律是从 0 和 1 开始向上滚动保存前两项，每个 fib(i) 只算一次”，说明这个参数组合可以命名成状态，算过之后后面直接复用。后面的优化只消掉这类重复工作，不能改变答案集合。

\textbf{优化条件。} 本题的关键观察是：自底向上只保存前两项，每次生成当前项后滚动更新。这句话必须能翻译成可维护状态；如果题目条件改变后观察不再成立，就不能继续套原来的写法。

\textbf{相邻状态。} 规律是从 0 和 1 开始向上滚动保存前两项，每个 fib(i) 只算一次。把这个特例继续拆开看：先问暴力搜索里哪些不同路径到达同一个后续问题，再给这个后续问题命名为状态。状态必须包含未来决策需要的全部信息，转移只从更小且已经算清的状态来。

\textbf{状态复用。} 把重复递归节点命名为状态；遍历顺序只是在保证依赖状态已经算完。本题落点是：规律是从 0 和 1 开始向上滚动保存前两项，每个 fib(i) 只算一次。手算时只改被新输入影响的那一小部分，而不是回头重算完整候选。

\textbf{指针和字段。} 状态字段要能说清：状态下标、状态值语义、初始化、转移来源、遍历顺序；空间压缩时还要指出读到的是旧值还是新值。手算时按这个协议跟踪：拿 n=5 手算，fib(5) 会递归调用 fib(4) 和 fib(3)，而 fib(4) 内部又会调用 fib(3)，同一个状态被重复计算。然后按协议复查：手算时先写一个很小的 DP 表或记忆化搜索记录。每个格子旁边写它来自哪些更小状态。

\textbf{代码落点。} 先写未压缩版本校验转移；压缩前后用小表对拍；不可达哨兵参与加法前要检查溢出。关键代码行必须能逐句翻译成“旧状态怎样变成新状态”，否则就只是模板。

\textbf{正确性。} 证明从最后一步开始：任意最优解的最后一步都在转移枚举中，转移来源已经由更小状态正确计算。

\textbf{边界实验。} 先用n 为零、n 为一、结果范围、递归爆栈做反例检查。实验时覆盖n 为零、n 为一、结果范围、递归爆栈，先打印完整 DP 表，再比较空间压缩版本。若压缩版不同，优先检查遍历顺序和旧值保存。

\textbf{复杂度变化。} 暴力递归会重复计算相同子问题，常见指数级；DP 把每个状态算一次，时间是状态数乘单个转移成本，空间是状态表大小或压缩后的大小。

\textbf{迁移追问。} 如果题目改成“若 n 极大，可用矩阵快速幂或快速倍增”，先判断原状态是否仍然覆盖新答案集合，再决定是微调边界、改 value 语义，还是换算法。

\noindent\textbf{LeetCode 516 最长回文子序列。}

\textbf{题目说明。} LeetCode 516 最长回文子序列 的题面入口要先还原成具体任务：子序列允许跳过字符，区间两端是否相等决定能否同时作为回文两端。

\textbf{样例入口。} 拿 s=bbbab 手算，两端 b 相等时可以同时作为回文两端，内部 bba 的最长回文再加 2；若两端不同，就只能丢左端或丢右端取较大。

\textbf{暴力基线。} 慢版本先写递归搜索树：节点表示处理位置、剩余资源和当前选择。只要不同路径反复到达同一个节点，就说明需要把这个节点命名成 DP 状态。放到本题就是：先按“子序列允许跳过字符，区间两端是否相等决定能否同时作为回文两端”完整试慢版本；样例里已经能看到“规律是 dp[left][right] 表示区间最长回文子序列，两端相等用内部区间加二，不等时取两个缩小区间的最大值”，所以优化前必须先能说清慢版本枚举了哪些候选。

\textbf{暴力过程。} 拿 s=bbbab 手算，两端 b 相等时可以同时作为回文两端，内部 bba 的最长回文再加 2；若两端不同，就只能丢左端或丢右端取较大。

\textbf{重复工作。} 题面模型：子序列允许跳过字符，区间两端是否相等决定能否同时作为回文两端。暴力版的慢点要落到本题：慢版本会在“子序列允许跳过字符，区间两端是否相等决定能否同时作为回文两端”的选择树里反复到达相同参数。样例规律是“规律是 dp[left][right] 表示区间最长回文子序列，两端相等用内部区间加二，不等时取两个缩小区间的最大值”，说明这个参数组合可以命名成状态，算过之后后面直接复用。后面的优化只消掉这类重复工作，不能改变答案集合。

\textbf{优化条件。} 本题的关键观察是：dp[left][right] 表示区间最长回文子序列，两端相等取内部加二，否则丢左或丢右取最大。这句话必须能翻译成可维护状态；如果题目条件改变后观察不再成立，就不能继续套原来的写法。

\textbf{相邻状态。} 规律是 dp[left][right] 表示区间最长回文子序列，两端相等用内部区间加二，不等时取两个缩小区间的最大值。把这个特例继续拆开看：先问暴力搜索里哪些不同路径到达同一个后续问题，再给这个后续问题命名为状态。状态必须包含未来决策需要的全部信息，转移只从更小且已经算清的状态来。

\textbf{状态复用。} 把重复递归节点命名为状态；遍历顺序只是在保证依赖状态已经算完。本题落点是：规律是 dp[left][right] 表示区间最长回文子序列，两端相等用内部区间加二，不等时取两个缩小区间的最大值。手算时只改被新输入影响的那一小部分，而不是回头重算完整候选。

\textbf{指针和字段。} 状态字段要能说清：状态下标、状态值语义、初始化、转移来源、遍历顺序；空间压缩时还要指出读到的是旧值还是新值。手算时按这个协议跟踪：拿 s=bbbab 手算，两端 b 相等时可以同时作为回文两端，内部 bba 的最长回文再加 2；若两端不同，就只能丢左端或丢右端取较大。然后按协议复查：手算时先写一个很小的 DP 表或记忆化搜索记录。每个格子旁边写它来自哪些更小状态。

\textbf{代码落点。} 先写未压缩版本校验转移；压缩前后用小表对拍；不可达哨兵参与加法前要检查溢出。关键代码行必须能逐句翻译成“旧状态怎样变成新状态”，否则就只是模板。

\textbf{正确性。} 证明从最后一步开始：任意最优解的最后一步都在转移枚举中，转移来源已经由更小状态正确计算。

\textbf{边界实验。} 先用空串、单字符、两端相等但内部为空、全相同字符、子串和子序列混淆做反例检查。实验时覆盖空串、单字符、两端相等但内部为空、全相同字符、子串和子序列混淆，先打印完整 DP 表，再比较空间压缩版本。若压缩版不同，优先检查遍历顺序和旧值保存。

\textbf{复杂度变化。} 暴力递归会重复计算相同子问题，常见指数级；DP 把每个状态算一次，时间是状态数乘单个转移成本，空间是状态表大小或压缩后的大小。

\textbf{迁移追问。} 如果题目改成“若求最长回文子串，状态必须保持连续性，不能用这个跳过模型”，先判断原状态是否仍然覆盖新答案集合，再决定是微调边界、改 value 语义，还是换算法。

\noindent\textbf{LeetCode 600 不含连续 1 的非负整数。}

\textbf{题目说明。} LeetCode 600 不含连续 1 的非负整数 的题面入口要先还原成具体任务：统计不超过 n 的二进制数，需要按位构造并记住上一位是否为一。

\textbf{样例入口。} 拿 n=5 的二进制 101 手算，不含连续 1 的数是 0、1、2、4、5。按高位扫描时，遇到某位为 1，可以先把这一位改成 0，后面任意填不含连续 1 的低位；若扫描中出现连续 1，就不能继续把 n 本身计入。

\textbf{暴力基线。} 慢版本先写递归搜索树：节点表示处理位置、剩余资源和当前选择。只要不同路径反复到达同一个节点，就说明需要把这个节点命名成 DP 状态。放到本题就是：先按“统计不超过 n 的二进制数，需要按位构造并记住上一位是否为一”完整试慢版本；样例里已经能看到“规律是数位 DP 需要保存上一位是否为 1 和是否贴着上界”，所以优化前必须先能说清慢版本枚举了哪些候选。

\textbf{暴力过程。} 拿 n=5 的二进制 101 手算，不含连续 1 的数是 0、1、2、4、5。按高位扫描时，遇到某位为 1，可以先把这一位改成 0，后面任意填不含连续 1 的低位；若扫描中出现连续 1，就不能继续把 n 本身计入。

\textbf{重复工作。} 题面模型：统计不超过 n 的二进制数，需要按位构造并记住上一位是否为一。暴力版的慢点要落到本题：慢版本会在“统计不超过 n 的二进制数，需要按位构造并记住上一位是否为一”的选择树里反复到达相同参数。样例规律是“规律是数位 DP 需要保存上一位是否为 1 和是否贴着上界”，说明这个参数组合可以命名成状态，算过之后后面直接复用。后面的优化只消掉这类重复工作，不能改变答案集合。

\textbf{优化条件。} 本题的关键观察是：数位 DP 或 Fibonacci 计数从高位扫描，遇到 1 时累加较小前缀数量，连续 1 时停止。这句话必须能翻译成可维护状态；如果题目条件改变后观察不再成立，就不能继续套原来的写法。

\textbf{相邻状态。} 规律是数位 DP 需要保存上一位是否为 1 和是否贴着上界。把这个特例继续拆开看：先问暴力搜索里哪些不同路径到达同一个后续问题，再给这个后续问题命名为状态。状态必须包含未来决策需要的全部信息，转移只从更小且已经算清的状态来。

\textbf{状态复用。} 把重复递归节点命名为状态；遍历顺序只是在保证依赖状态已经算完。本题落点是：规律是数位 DP 需要保存上一位是否为 1 和是否贴着上界。手算时只改被新输入影响的那一小部分，而不是回头重算完整候选。

\textbf{指针和字段。} 状态字段要能说清：状态下标、状态值语义、初始化、转移来源、遍历顺序；空间压缩时还要指出读到的是旧值还是新值。手算时按这个协议跟踪：拿 n=5 的二进制 101 手算，不含连续 1 的数是 0、1、2、4、5。按高位扫描时，遇到某位为 1，可以先把这一位改成 0，后面任意填不含连续 1 的低位；若扫描中出现连续 1，就不能继续把 n 本身计入。然后按协议复查：手算时先写一个很小的 DP 表或记忆化搜索记录。每个格子旁边写它来自哪些更小状态。

\textbf{代码落点。} 先写未压缩版本校验转移；压缩前后用小表对拍；不可达哨兵参与加法前要检查溢出。关键代码行必须能逐句翻译成“旧状态怎样变成新状态”，否则就只是模板。

\textbf{正确性。} 证明从最后一步开始：任意最优解的最后一步都在转移枚举中，转移来源已经由更小状态正确计算。

\textbf{边界实验。} 先用n 为零、最高位、连续一提前停止、包含 n 本身、整数位宽做反例检查。实验时覆盖n 为零、最高位、连续一提前停止、包含 n 本身、整数位宽，先打印完整 DP 表，再比较空间压缩版本。若压缩版不同，优先检查遍历顺序和旧值保存。

\textbf{复杂度变化。} 暴力递归会重复计算相同子问题，常见指数级；DP 把每个状态算一次，时间是状态数乘单个转移成本，空间是状态表大小或压缩后的大小。

\textbf{迁移追问。} 如果题目改成“若禁止其他数字模式，状态要记录更长后缀或自动机”，先判断原状态是否仍然覆盖新答案集合，再决定是微调边界、改 value 语义，还是换算法。

\noindent\textbf{LeetCode 714 买卖股票含手续费。}

\textbf{题目说明。} LeetCode 714 买卖股票含手续费 的题面入口要先还原成具体任务：手续费改变交易收益，仍可用持有和不持有状态。

\textbf{样例入口。} 拿 [1, 3, 2, 8, 4, 9]、fee=2 手算，8 卖出时净赚 5；手续费让频繁小波动不一定值得交易。

\textbf{暴力基线。} 慢版本先写递归搜索树：节点表示处理位置、剩余资源和当前选择。只要不同路径反复到达同一个节点，就说明需要把这个节点命名成 DP 状态。放到本题就是：先按“手续费改变交易收益，仍可用持有和不持有状态”完整试慢版本；样例里已经能看到“规律是每天维护 hold 和 cash，卖出时扣手续费，买入时从 cash 转到 hold”，所以优化前必须先能说清慢版本枚举了哪些候选。

\textbf{暴力过程。} 拿 [1, 3, 2, 8, 4, 9]、fee=2 手算，8 卖出时净赚 5；手续费让频繁小波动不一定值得交易。

\textbf{重复工作。} 题面模型：手续费改变交易收益，仍可用持有和不持有状态。暴力版的慢点要落到本题：慢版本会在“手续费改变交易收益，仍可用持有和不持有状态”的选择树里反复到达相同参数。样例规律是“规律是每天维护 hold 和 cash，卖出时扣手续费，买入时从 cash 转到 hold”，说明这个参数组合可以命名成状态，算过之后后面直接复用。后面的优化只消掉这类重复工作，不能改变答案集合。

\textbf{优化条件。} 本题的关键观察是：卖出时扣费或买入时扣费都可以，但不要重复扣。这句话必须能翻译成可维护状态；如果题目条件改变后观察不再成立，就不能继续套原来的写法。

\textbf{相邻状态。} 规律是每天维护 hold 和 cash，卖出时扣手续费，买入时从 cash 转到 hold。把这个特例继续拆开看：先问暴力搜索里哪些不同路径到达同一个后续问题，再给这个后续问题命名为状态。状态必须包含未来决策需要的全部信息，转移只从更小且已经算清的状态来。

\textbf{状态复用。} 把重复递归节点命名为状态；遍历顺序只是在保证依赖状态已经算完。本题落点是：规律是每天维护 hold 和 cash，卖出时扣手续费，买入时从 cash 转到 hold。手算时只改被新输入影响的那一小部分，而不是回头重算完整候选。

\textbf{指针和字段。} 状态字段要能说清：状态下标、状态值语义、初始化、转移来源、遍历顺序；空间压缩时还要指出读到的是旧值还是新值。手算时按这个协议跟踪：拿 [1, 3, 2, 8, 4, 9]、fee=2 手算，8 卖出时净赚 5；手续费让频繁小波动不一定值得交易。然后按协议复查：手算时先写一个很小的 DP 表或记忆化搜索记录。每个格子旁边写它来自哪些更小状态。

\textbf{代码落点。} 先写未压缩版本校验转移；压缩前后用小表对拍；不可达哨兵参与加法前要检查溢出。关键代码行必须能逐句翻译成“旧状态怎样变成新状态”，否则就只是模板。

\textbf{正确性。} 证明从最后一步开始：任意最优解的最后一步都在转移枚举中，转移来源已经由更小状态正确计算。

\textbf{边界实验。} 先用手续费为零、价格震荡、小利润不足手续费做反例检查。实验时覆盖手续费为零、价格震荡、小利润不足手续费，先打印完整 DP 表，再比较空间压缩版本。若压缩版不同，优先检查遍历顺序和旧值保存。

\textbf{复杂度变化。} 暴力递归会重复计算相同子问题，常见指数级；DP 把每个状态算一次，时间是状态数乘单个转移成本，空间是状态表大小或压缩后的大小。

\textbf{迁移追问。} 如果题目改成“这是状态机比局部贪心更稳的例子”，先判断原状态是否仍然覆盖新答案集合，再决定是微调边界、改 value 语义，还是换算法。

\noindent\textbf{LeetCode 790 多米诺和托米诺平铺。}

\textbf{题目说明。} LeetCode 790 多米诺和托米诺平铺 的题面入口要先还原成具体任务：铺满前 i 列时，边界可能完整，也可能上缺或下缺一个格子。

\textbf{样例入口。} 拿 n=3 手算，只有多米诺可以铺出若干完整状态，加入 L 形托米诺后会产生上缺一格或下缺一格的半填状态。

\textbf{暴力基线。} 慢版本先写递归搜索树：节点表示处理位置、剩余资源和当前选择。只要不同路径反复到达同一个节点，就说明需要把这个节点命名成 DP 状态。放到本题就是：先按“铺满前 i 列时，边界可能完整，也可能上缺或下缺一个格子”完整试慢版本；样例里已经能看到“规律是 DP 不能只保存完整列，还要保存缺口状态；完整状态和两个缺口状态互相转移”，所以优化前必须先能说清慢版本枚举了哪些候选。

\textbf{暴力过程。} 拿 n=3 手算，只有多米诺可以铺出若干完整状态，加入 L 形托米诺后会产生上缺一格或下缺一格的半填状态。

\textbf{重复工作。} 题面模型：铺满前 i 列时，边界可能完整，也可能上缺或下缺一个格子。暴力版的慢点要落到本题：慢版本会在“铺满前 i 列时，边界可能完整，也可能上缺或下缺一个格子”的选择树里反复到达相同参数。样例规律是“规律是 DP 不能只保存完整列，还要保存缺口状态；完整状态和两个缺口状态互相转移”，说明这个参数组合可以命名成状态，算过之后后面直接复用。后面的优化只消掉这类重复工作，不能改变答案集合。

\textbf{优化条件。} 本题的关键观察是：维护完整状态和缺口状态，按最后一块多米诺或托米诺转移。这句话必须能翻译成可维护状态；如果题目条件改变后观察不再成立，就不能继续套原来的写法。

\textbf{相邻状态。} 规律是 DP 不能只保存完整列，还要保存缺口状态；完整状态和两个缺口状态互相转移。把这个特例继续拆开看：先问暴力搜索里哪些不同路径到达同一个后续问题，再给这个后续问题命名为状态。状态必须包含未来决策需要的全部信息，转移只从更小且已经算清的状态来。

\textbf{状态复用。} 把重复递归节点命名为状态；遍历顺序只是在保证依赖状态已经算完。本题落点是：规律是 DP 不能只保存完整列，还要保存缺口状态；完整状态和两个缺口状态互相转移。手算时只改被新输入影响的那一小部分，而不是回头重算完整候选。

\textbf{指针和字段。} 状态字段要能说清：状态下标、状态值语义、初始化、转移来源、遍历顺序；空间压缩时还要指出读到的是旧值还是新值。手算时按这个协议跟踪：拿 n=3 手算，只有多米诺可以铺出若干完整状态，加入 L 形托米诺后会产生上缺一格或下缺一格的半填状态。然后按协议复查：手算时先写一个很小的 DP 表或记忆化搜索记录。每个格子旁边写它来自哪些更小状态。

\textbf{代码落点。} 先写未压缩版本校验转移；压缩前后用小表对拍；不可达哨兵参与加法前要检查溢出。关键代码行必须能逐句翻译成“旧状态怎样变成新状态”，否则就只是模板。

\textbf{正确性。} 证明从最后一步开始：任意最优解的最后一步都在转移枚举中，转移来源已经由更小状态正确计算。

\textbf{边界实验。} 先用n 为一、n 为二、取模、缺口上下对称、初始化做反例检查。实验时覆盖n 为一、n 为二、取模、缺口上下对称、初始化，先打印完整 DP 表，再比较空间压缩版本。若压缩版不同，优先检查遍历顺序和旧值保存。

\textbf{复杂度变化。} 暴力递归会重复计算相同子问题，常见指数级；DP 把每个状态算一次，时间是状态数乘单个转移成本，空间是状态表大小或压缩后的大小。

\textbf{迁移追问。} 如果题目改成“更多形状的铺砖题通常要改成轮廓状态压缩 DP”，先判断原状态是否仍然覆盖新答案集合，再决定是微调边界、改 value 语义，还是换算法。

\noindent\textbf{LeetCode 1143 最长公共子序列。}

\textbf{题目说明。} LeetCode 1143 最长公共子序列 的题面入口要先还原成具体任务：状态表示两个前缀的最长公共子序列长度。

\textbf{样例入口。} 拿 abcde 和 ace 手算，末尾 e 相等时答案来自两个前缀 abcd 和 ac 加一；若末尾不同，就取删掉一侧末尾后的较大值。

\textbf{暴力基线。} 暴力递归枚举两个字符串的子序列选择，或从末尾开始相等就选、不等就分支删一边。

\textbf{暴力过程。} 以 text1=abcde、text2=ace 为例，末尾 e 相等时，答案来自 abcd 和 ac 的 LCS 加一；若末尾不同，就只能丢掉一侧末尾试更优。

\textbf{重复工作。} 题面模型：状态表示两个前缀的最长公共子序列长度。暴力版的慢点要落到本题：浪费在很多前缀对会被重复计算。只要两个前缀长度 i、j 固定，它们的 LCS 长度就是固定状态。后面的优化只消掉这类重复工作，不能改变答案集合。

\textbf{优化条件。} 本题的关键观察是：末尾相同取左上加一，不同取上方和左方最大。这句话必须能翻译成可维护状态；如果题目条件改变后观察不再成立，就不能继续套原来的写法。

\textbf{相邻状态。} 规律是 dp[i][j] 表示两个前缀的 LCS 长度，转移只看左、上、左上。把这个特例继续拆开看：先问暴力搜索里哪些不同路径到达同一个后续问题，再给这个后续问题命名为状态。状态必须包含未来决策需要的全部信息，转移只从更小且已经算清的状态来。

\textbf{状态复用。} dp[i][j] 表示 text1 前 i 个字符和 text2 前 j 个字符的 LCS 长度。相等看左上加一，不等取上方和左方最大。手算时只改被新输入影响的那一小部分，而不是回头重算完整候选。

\textbf{指针和字段。} 状态字段要能说清：i/j 是前缀长度，dp[i-1][j] 表示丢 text1 末尾，dp[i][j-1] 表示丢 text2 末尾，dp[i-1][j-1] 表示两个末尾配对。手算时按这个协议跟踪：a 与 a 相等得到 1；到 abcde 和 ace 的最后，e 与 e 相等，所以由 abcd/ac 的答案加一。

\textbf{代码落点。} 关键是表多开一行一列表示空前缀，循环从 1 开始，访问字符用 text[i-1]。关键代码行必须能逐句翻译成“旧状态怎样变成新状态”，否则就只是模板。

\textbf{正确性。} 任意 LCS 对末尾字符只有配对或至少丢掉一边两类；转移覆盖这两类并取最大。

\textbf{边界实验。} 先用空串、完全相同、无公共字符、子序列不是子串做反例检查。实验时覆盖空串、完全相同、无公共字符、子序列不是子串，先打印完整 DP 表，再比较空间压缩版本。若压缩版不同，优先检查遍历顺序和旧值保存。

\textbf{复杂度变化。} 暴力指数级；DP O(mn) 时间和空间，可滚动到 O(min(m, n)) 空间。

\textbf{迁移追问。} 如果题目改成“若要求还原序列，需要从表右下角反向追踪选择”，先判断原状态是否仍然覆盖新答案集合，再决定是微调边界、改 value 语义，还是换算法。

\noindent\textbf{LeetCode 1235 规划兼职工作。}

\textbf{题目说明。} LeetCode 1235 规划兼职工作 的题面入口要先还原成具体任务：每份工作有开始、结束和收益，选择不冲突工作使收益最大。

\textbf{样例入口。} 拿 jobs (1, 3, 50), (2, 4, 10), (3, 5, 40), (3, 6, 70) 手算，选 (1, 3, 50) 后还能接 (3, 6, 70)，总 120。

\textbf{暴力基线。} 慢版本先写递归搜索树：节点表示处理位置、剩余资源和当前选择。只要不同路径反复到达同一个节点，就说明需要把这个节点命名成 DP 状态。放到本题就是：先按“每份工作有开始、结束和收益，选择不冲突工作使收益最大”完整试慢版本；样例里已经能看到“规律是按结束时间排序，dp[i] 在不选当前和选当前加上前一个不冲突工作之间取最大”，所以优化前必须先能说清慢版本枚举了哪些候选。

\textbf{暴力过程。} 拿 jobs (1, 3, 50), (2, 4, 10), (3, 5, 40), (3, 6, 70) 手算，选 (1, 3, 50) 后还能接 (3, 6, 70)，总 120。

\textbf{重复工作。} 题面模型：每份工作有开始、结束和收益，选择不冲突工作使收益最大。暴力版的慢点要落到本题：慢版本会在“每份工作有开始、结束和收益，选择不冲突工作使收益最大”的选择树里反复到达相同参数。样例规律是“规律是按结束时间排序，dp[i] 在不选当前和选当前加上前一个不冲突工作之间取最大”，说明这个参数组合可以命名成状态，算过之后后面直接复用。后面的优化只消掉这类重复工作，不能改变答案集合。

\textbf{优化条件。} 本题的关键观察是：按结束时间排序，dp[i] 在不选当前和选当前加上兼容前驱之间取最大，前驱用二分找。这句话必须能翻译成可维护状态；如果题目条件改变后观察不再成立，就不能继续套原来的写法。

\textbf{相邻状态。} 规律是按结束时间排序，dp[i] 在不选当前和选当前加上前一个不冲突工作之间取最大。把这个特例继续拆开看：先问暴力搜索里哪些不同路径到达同一个后续问题，再给这个后续问题命名为状态。状态必须包含未来决策需要的全部信息，转移只从更小且已经算清的状态来。

\textbf{状态复用。} 把重复递归节点命名为状态；遍历顺序只是在保证依赖状态已经算完。本题落点是：规律是按结束时间排序，dp[i] 在不选当前和选当前加上前一个不冲突工作之间取最大。手算时只改被新输入影响的那一小部分，而不是回头重算完整候选。

\textbf{指针和字段。} 状态字段要能说清：状态下标、状态值语义、初始化、转移来源、遍历顺序；空间压缩时还要指出读到的是旧值还是新值。手算时按这个协议跟踪：拿 jobs (1, 3, 50), (2, 4, 10), (3, 5, 40), (3, 6, 70) 手算，选 (1, 3, 50) 后还能接 (3, 6, 70)，总 120。然后按协议复查：手算时先写一个很小的 DP 表或记忆化搜索记录。每个格子旁边写它来自哪些更小状态。

\textbf{代码落点。} 先写未压缩版本校验转移；压缩前后用小表对拍；不可达哨兵参与加法前要检查溢出。关键代码行必须能逐句翻译成“旧状态怎样变成新状态”，否则就只是模板。

\textbf{正确性。} 证明从最后一步开始：任意最优解的最后一步都在转移枚举中，转移来源已经由更小状态正确计算。

\textbf{边界实验。} 先用结束等于开始是否兼容、收益相同、工作排序、空列表做反例检查。实验时覆盖结束等于开始是否兼容、收益相同、工作排序、空列表，先打印完整 DP 表，再比较空间压缩版本。若压缩版不同，优先检查遍历顺序和旧值保存。

\textbf{复杂度变化。} 暴力递归会重复计算相同子问题，常见指数级；DP 把每个状态算一次，时间是状态数乘单个转移成本，空间是状态表大小或压缩后的大小。

\textbf{迁移追问。} 如果题目改成“这是排序、二分和 DP 的组合，不是按收益贪心”，先判断原状态是否仍然覆盖新答案集合，再决定是微调边界、改 value 语义，还是换算法。

\noindent\textbf{LeetCode 1425 带限制的子序列和。}

\textbf{题目说明。} LeetCode 1425 带限制的子序列和 的题面入口要先还原成具体任务：dp[i] 只依赖前 k 个位置的最大 dp 值，转移范围是滑动窗口。

\textbf{样例入口。} 拿 nums=[10, 2, -10, 5, 20]、k=2 手算，dp[3] 的前驱只能在下标 1 和 2 中选，dp[4] 又能接 dp[3] 得到 37。暴力每个 i 回看前 k 个 dp；单调队列保留窗口内 dp 最大的下标。

\textbf{暴力基线。} 慢版本先写递归搜索树：节点表示处理位置、剩余资源和当前选择。只要不同路径反复到达同一个节点，就说明需要把这个节点命名成 DP 状态。放到本题就是：先按“dp[i] 只依赖前 k 个位置的最大 dp 值，转移范围是滑动窗口”完整试慢版本；样例里已经能看到“规律是转移 dp[i]=nums[i]+max(0, max dp[j])，其中 j 在 i-k 到 i-1 内”，所以优化前必须先能说清慢版本枚举了哪些候选。

\textbf{暴力过程。} 拿 nums=[10, 2, -10, 5, 20]、k=2 手算，dp[3] 的前驱只能在下标 1 和 2 中选，dp[4] 又能接 dp[3] 得到 37。暴力每个 i 回看前 k 个 dp；单调队列保留窗口内 dp 最大的下标。

\textbf{重复工作。} 题面模型：dp[i] 只依赖前 k 个位置的最大 dp 值，转移范围是滑动窗口。暴力版的慢点要落到本题：慢版本会在“dp[i] 只依赖前 k 个位置的最大 dp 值，转移范围是滑动窗口”的选择树里反复到达相同参数。样例规律是“规律是转移 dp[i]=nums[i]+max(0, max dp[j])，其中 j 在 i-k 到 i-1 内”，说明这个参数组合可以命名成状态，算过之后后面直接复用。后面的优化只消掉这类重复工作，不能改变答案集合。

\textbf{优化条件。} 本题的关键观察是：单调队列保存候选 dp 下标，过期弹队首，被更大 dp 支配的候选从队尾删除。这句话必须能翻译成可维护状态；如果题目条件改变后观察不再成立，就不能继续套原来的写法。

\textbf{相邻状态。} 规律是转移 dp[i]=nums[i]+max(0, max dp[j])，其中 j 在 i-k 到 i-1 内。把这个特例继续拆开看：先问暴力搜索里哪些不同路径到达同一个后续问题，再给这个后续问题命名为状态。状态必须包含未来决策需要的全部信息，转移只从更小且已经算清的状态来。

\textbf{状态复用。} 把重复递归节点命名为状态；遍历顺序只是在保证依赖状态已经算完。本题落点是：规律是转移 dp[i]=nums[i]+max(0, max dp[j])，其中 j 在 i-k 到 i-1 内。手算时只改被新输入影响的那一小部分，而不是回头重算完整候选。

\textbf{指针和字段。} 状态字段要能说清：状态下标、状态值语义、初始化、转移来源、遍历顺序；空间压缩时还要指出读到的是旧值还是新值。手算时按这个协议跟踪：拿 nums=[10, 2, -10, 5, 20]、k=2 手算，dp[3] 的前驱只能在下标 1 和 2 中选，dp[4] 又能接 dp[3] 得到 37。暴力每个 i 回看前 k 个 dp；单调队列保留窗口内 dp 最大的下标。然后按协议复查：手算时先写一个很小的 DP 表或记忆化搜索记录。每个格子旁边写它来自哪些更小状态。

\textbf{代码落点。} 先写未压缩版本校验转移；压缩前后用小表对拍；不可达哨兵参与加法前要检查溢出。关键代码行必须能逐句翻译成“旧状态怎样变成新状态”，否则就只是模板。

\textbf{正确性。} 证明从最后一步开始：任意最优解的最后一步都在转移枚举中，转移来源已经由更小状态正确计算。

\textbf{边界实验。} 先用全负数、k 为一、k 大于数组长度、答案必须非空、队列过期顺序做反例检查。实验时覆盖全负数、k 为一、k 大于数组长度、答案必须非空、队列过期顺序，先打印完整 DP 表，再比较空间压缩版本。若压缩版不同，优先检查遍历顺序和旧值保存。

\textbf{复杂度变化。} 暴力递归会重复计算相同子问题，常见指数级；DP 把每个状态算一次，时间是状态数乘单个转移成本，空间是状态表大小或压缩后的大小。

\textbf{迁移追问。} 如果题目改成“若窗口内需要最小值或其他可合并统计，队列维护规则要重证”，先判断原状态是否仍然覆盖新答案集合，再决定是微调边界、改 value 语义，还是换算法。

\subsection{实验复盘}

追问通常会问状态是否足够、初始化为什么这样、空间压缩读到的是旧值还是新值。请准备从暴力搜索树指出重复子问题。

复盘本节时先从 LeetCode 10 正则表达式匹配、LeetCode 44 通配符匹配、LeetCode 53 最大子数组和 开始：每道题都先手写一个能覆盖全部候选的暴力版，再指出优化结构具体保存了哪一段历史信息，最后用相邻案例比较为什么不能互换解法。

\topic{背包模型}
这一节先看 LeetCode 279 完全平方数、LeetCode 322 零钱兑换、LeetCode 377 组合总和 IV 这几道 LeetCode 题，再整理同一题型的分流和推导。阅读顺序不是先背原理，而是先看案例的暴力瓶颈，再把重复出现的观察抽象成方法。

\subsection{原理和适用条件}

以 LeetCode 279 完全平方数、LeetCode 322 零钱兑换、LeetCode 377 组合总和 IV 为主线：LeetCode 279 完全平方数：模型是“完全平方数可以看成无限使用的物品，目标是凑出 n 的最少数量。”，优化抓手是“完全背包最少物品数，或把数字看成图节点做 BFS。”；LeetCode 322 零钱兑换：模型是“金额是容量，硬币可无限使用，目标是最少硬币数。”，优化抓手是“dp[value] 从 value 减 coin 的已知状态转移。”；LeetCode 377 组合总和 IV：模型是“题目统计排列数，选择顺序不同算不同方案。”，优化抓手是“外层遍历容量，内层遍历数字，让每个容量的方案按最后一个数累加。”。这些案例共同推出的规律是：先枚举物品选择，再把剩余容量作为状态。0-1 背包倒序容量，完全背包正序容量，组合和排列的循环顺序不能混用。 方法名只是复盘后的标签，真正要掌握的是从题面约束、暴力失败、状态设计到边界反例的推导链。

\begin{principlebox}{图示：从 LeetCode 案例到 背包模型推导链}
\begin{tabular}{@{}p{0.18\linewidth}p{0.74\linewidth}@{}}
暴力基线 & 枚举候选、完整模拟或逐个检查，先保证覆盖全部可能答案。\\
瓶颈定位 & 瓶颈是阶段相同、剩余容量相同的子问题反复出现。容量把未来选择压缩成一个数字，因此可以用表保存。\\
结构选择 & 0-1 背包每个物品只能用一次，一维压缩时容量倒序；完全背包物品可无限使用，容量正序。计数组合和计数排列的循环顺序不同。\\
不变量 & 处理到第 i 个物品后，\code{dp[cap]} 表示只使用前 i 类物品能达到的最优值或方案数。倒序是为了不在同一轮重复使用当前物品。\\
代码落地 & 容量通常与数值大小有关，所以复杂度是伪多项式。实现时先判断总和奇偶、目标范围和负数是否存在。\\
\end{tabular}
\end{principlebox}

\subsection{LeetCode 案例分流}

\begin{principlebox}{从 LeetCode 案例分流：背包模型}
\textbf{0-1 背包可达性。}
代表案例：LeetCode 416 分割等和子集、LeetCode 494 目标和。先看这些题的题面条件：每个物品最多使用一次，问能否达到目标容量。由此推出的处理方式是：一维压缩时容量倒序。风险点：容量正序会把同一物品用多次。

\textbf{完全背包最优值。}
代表案例：LeetCode 322 零钱兑换、LeetCode 279 完全平方数。先看这些题的题面条件：物品可以无限使用，问最少数量或最大价值。由此推出的处理方式是：容量正序，让当前物品可在同一轮继续使用。风险点：不可达哨兵要防溢出。

\textbf{完全背包计数。}
代表案例：LeetCode 518 零钱兑换 II、LeetCode 377 组合总和 IV。先看这些题的题面条件：物品可以无限使用，问组合数或排列数。由此推出的处理方式是：组合数外层物品，排列数外层容量。风险点：循环顺序直接改变题意。

\textbf{二维容量和约束扩展。}
代表案例：LeetCode 474 一和零、LeetCode 1049 最后一块石头重量 II。先看这些题的题面条件：物品消耗不止一种资源，或目标需要多个维度。由此推出的处理方式是：状态增加容量维度或压缩为二维表。风险点：维度增加后复杂度和容量乘积相关。

\end{principlebox}

\subsection{案例矩阵}

本节案例覆盖 LeetCode 279 完全平方数、LeetCode 322 零钱兑换、LeetCode 377 组合总和 IV、LeetCode 416 分割等和子集、LeetCode 474 一和零、LeetCode 494 目标和、LeetCode 518 零钱兑换 II、LeetCode 879 盈利计划、LeetCode 956 最高的广告牌、LeetCode 1049 最后一块石头的重量 II、LeetCode 1155 掷骰子等于目标和的方法数、LeetCode 1449 数位成本和为目标值的最大数字。阅读时不要按题号背答案，而要先判断题面落在哪个分流场景：查询目标是什么，输入约束给了什么单调性或等价关系，暴力慢在哪里，优化结构保存了什么状态。

\begin{principlebox}{案例矩阵}
\textbf{LeetCode 279 完全平方数。}
完全平方数可以看成无限使用的物品，目标是凑出 n 的最少数量。单点风险：n 为零或一、完全平方数本身、较大 n。

\textbf{LeetCode 322 零钱兑换。}
金额是容量，硬币可无限使用，目标是最少硬币数。单点风险：无法凑出、amount 为零、硬币面额大于目标、哨兵溢出。

\textbf{LeetCode 377 组合总和 IV。}
题目统计排列数，选择顺序不同算不同方案。单点风险：target 为零、数字大于目标、方案数溢出、循环顺序。

\textbf{LeetCode 416 分割等和子集。}
总和一半是背包容量，每个数只能使用一次。单点风险：总和为奇数、目标为零、数值较大、复杂度依赖 target。

\textbf{LeetCode 474 一和零。}
每个字符串消耗零和一两种容量，每个字符串最多选一次。单点风险：字符串全零或全一、容量为零、倒序方向、计数预处理。

\textbf{LeetCode 494 目标和。}
正负号选择可代数转换为选若干数凑出特定正和。单点风险：零元素会让方案数翻倍、目标为负、容量为零。

\textbf{LeetCode 518 零钱兑换 II。}
硬币无限使用，问组合数而不是最少数量。单点风险：amount 为零、无硬币、组合数较大、循环顺序反例。

\textbf{LeetCode 879 盈利计划。}
每个项目消耗人数并贡献利润，利润达到下限后可截断。单点风险：minProfit 为零、人数不足、方案数取模、项目只能选一次。

\textbf{LeetCode 956 最高的广告牌。}
两组钢筋高度相等时得到广告牌，高度差是关键状态。单点风险：空数组、无法平衡、重复长度、差值更新顺序。

\textbf{LeetCode 1049 最后一块石头的重量 II。}
把石头分成两堆，最终差值等于两堆重量差。单点风险：单块石头、总和为奇偶、多个相同重量、容量方向。

\textbf{LeetCode 1155 掷骰子等于目标和的方法数。}
多个骰子每个只能取一到 faces 的点数，问目标和方案数。单点风险：target 太小或太大、取模、骰子数为一、滚动数组方向。

\textbf{LeetCode 1449 数位成本和为目标值的最大数字。}
每个数字可重复选择，容量是成本，目标是组成字典序和长度最大的数字。单点风险：无法组成目标、多个数字同成本、target 为零、字符串比较。

\end{principlebox}

\subsection{共性落点}

\begin{principlebox}{本组共性落点}
\textbf{慢版本。} 暴力枚举每个物品选或不选，或者枚举每种硬币使用次数。物品数稍大时，搜索树会迅速膨胀。
\textbf{瓶颈。} 瓶颈是阶段相同、剩余容量相同的子问题反复出现。容量把未来选择压缩成一个数字，因此可以用表保存。
\textbf{状态字段。} 物品阶段、容量、状态值、循环方向、取模或不可达哨兵；组合和排列必须用不同循环顺序表达。
\textbf{不变量。} 处理到第 i 个物品后，\code{dp[cap]} 表示只使用前 i 类物品能达到的最优值或方案数。倒序是为了不在同一轮重复使用当前物品。
\textbf{实现检查。} 0-1 倒序、完全正序、排列外层容量、组合外层物品；每种循环方向都要能用一个小例子解释。
\textbf{C++ 落地。} 容量通常与数值大小有关，所以复杂度是伪多项式。实现时先判断总和奇偶、目标范围和负数是否存在。
\textbf{证明抓手。} 证明从当前物品使用次数开始：0-1、完全、计数和最值的循环语义必须和题意一致。
\end{principlebox}

\subsection{案例推导}

\noindent\textbf{LeetCode 279 完全平方数。}

\textbf{题目说明。} LeetCode 279 完全平方数 的题面入口要先还原成具体任务：完全平方数可以看成无限使用的物品，目标是凑出 n 的最少数量。

\textbf{样例入口。} 拿 n=12 手算，12 可以由 4+4+4 组成，不能只贪心选 9 后剩 3。

\textbf{暴力基线。} 慢版本枚举每个物品选不选、选几次，并记录阶段和剩余容量。相同阶段和容量反复出现时，就可以把指数搜索压成表。放到本题就是：先按“完全平方数可以看成无限使用的物品，目标是凑出 n 的最少数量”完整试慢版本；样例里已经能看到“规律是完全背包，dp[x] 表示和为 x 的最少平方数，枚举每个平方数时允许重复使用”，所以优化前必须先能说清慢版本枚举了哪些候选。

\textbf{暴力过程。} 拿 n=12 手算，12 可以由 4+4+4 组成，不能只贪心选 9 后剩 3。

\textbf{重复工作。} 题面模型：完全平方数可以看成无限使用的物品，目标是凑出 n 的最少数量。暴力版的慢点要落到本题：慢版本会围绕“完全平方数可以看成无限使用的物品，目标是凑出 n 的最少数量”枚举物品使用情况。样例规律是“规律是完全背包，dp[x] 表示和为 x 的最少平方数，枚举每个平方数时允许重复使用”，说明阶段、容量和使用次数已经足够描述后续选择，没必要反复展开同一剩余容量。后面的优化只消掉这类重复工作，不能改变答案集合。

\textbf{优化条件。} 本题的关键观察是：完全背包最少物品数，或把数字看成图节点做 BFS。这句话必须能翻译成可维护状态；如果题目条件改变后观察不再成立，就不能继续套原来的写法。

\textbf{相邻状态。} 规律是完全背包，dp[x] 表示和为 x 的最少平方数，枚举每个平方数时允许重复使用。把这个特例继续拆开看：阶段是物品，容量是资源，状态值是可达、最值还是计数。一个物品能否在同一轮再次使用，直接决定容量循环方向；组合和排列也由循环顺序表达。

\textbf{状态复用。} 把选择树压成阶段和容量；物品是否可重复使用决定一维表的遍历方向。本题落点是：规律是完全背包，dp[x] 表示和为 x 的最少平方数，枚举每个平方数时允许重复使用。手算时只改被新输入影响的那一小部分，而不是回头重算完整候选。

\textbf{指针和字段。} 状态字段要能说清：物品阶段、容量、状态值、循环方向、取模或不可达哨兵；组合和排列必须用不同循环顺序表达。手算时按这个协议跟踪：拿 n=12 手算，12 可以由 4+4+4 组成，不能只贪心选 9 后剩 3。然后按协议复查：手算时用一个容量很小的例子比较正序和倒序。只要循环方向改变后同一物品能否重复使用发生变化，就说明方向不是细节。

\textbf{代码落点。} 0-1 倒序、完全正序、排列外层容量、组合外层物品；每种循环方向都要能用一个小例子解释。关键代码行必须能逐句翻译成“旧状态怎样变成新状态”，否则就只是模板。

\textbf{正确性。} 证明从当前物品使用次数开始：0-1、完全、计数和最值的循环语义必须和题意一致。

\textbf{边界实验。} 先用n 为零或一、完全平方数本身、较大 n做反例检查。实验时覆盖n 为零或一、完全平方数本身、较大 n，并故意把容量方向写反构造最小反例。这样能确认循环方向来自题意，不是背模板。

\textbf{复杂度变化。} 暴力枚举物品选择指数级；背包 DP 通常按物品数乘容量或多维容量乘积计算，这是伪多项式复杂度。

\textbf{迁移追问。} 如果题目改成“数学四平方定理可给更强解法，但面试 DP 更通用”，先判断原状态是否仍然覆盖新答案集合，再决定是微调边界、改 value 语义，还是换算法。

\noindent\textbf{LeetCode 322 零钱兑换。}

\textbf{题目说明。} LeetCode 322 零钱兑换 的题面入口要先还原成具体任务：金额是容量，硬币可无限使用，目标是最少硬币数。

\textbf{样例入口。} 拿 coins=[1, 2, 5]、amount=11 手算，最后一步可以用 1、2 或 5，若用 5，则来源是 dp[6]+1。

\textbf{暴力基线。} 暴力递归枚举最后使用哪枚硬币，或者枚举每种硬币使用多少次，直到金额变成 0。

\textbf{暴力过程。} 以 coins=[1, 2, 5]、amount=11 为例，若最后一枚是 5，前面必须先凑出 6；若最后一枚是 2，前面必须先凑出 9。

\textbf{重复工作。} 题面模型：金额是容量，硬币可无限使用，目标是最少硬币数。暴力版的慢点要落到本题：浪费在同一个剩余金额会从很多硬币序列到达。凑出金额 x 的最少硬币数和之前路径无关，只和 x 有关。后面的优化只消掉这类重复工作，不能改变答案集合。

\textbf{优化条件。} 本题的关键观察是：dp[value] 从 value 减 coin 的已知状态转移。这句话必须能翻译成可维护状态；如果题目条件改变后观察不再成立，就不能继续套原来的写法。

\textbf{相邻状态。} 规律是 dp[x] 表示凑成 x 的最少硬币数，完全背包允许同一硬币重复使用；不可达状态不能参与加一。把这个特例继续拆开看：阶段是物品，容量是资源，状态值是可达、最值还是计数。一个物品能否在同一轮再次使用，直接决定容量循环方向；组合和排列也由循环顺序表达。

\textbf{状态复用。} dp[x] 表示凑出金额 x 的最少硬币数，初始化 dp[0]=0、其他为 INF。对每个金额尝试所有 coin。手算时只改被新输入影响的那一小部分，而不是回头重算完整候选。

\textbf{指针和字段。} 状态字段要能说清：x 是当前金额，coin 是最后一枚硬币，dp[x-coin]+1 是使用 coin 的候选答案。手算时按这个协议跟踪：dp[5] 可由 coin=5 得到 1；dp[11] 可由 dp[6]+1、dp[9]+1、dp[10]+1 取最小。

\textbf{代码落点。} 关键是不可达 INF 不能直接加一造成溢出；完全背包一维写法可金额正序，也可金额外层逐项求最小。关键代码行必须能逐句翻译成“旧状态怎样变成新状态”，否则就只是模板。

\textbf{正确性。} 任意最优方案都有一枚最后硬币 coin，去掉它后剩余金额 x-coin 必须也是最优子问题，否则可替换变更优。

\textbf{边界实验。} 先用无法凑出、amount 为零、硬币面额大于目标、哨兵溢出做反例检查。实验时覆盖无法凑出、amount 为零、硬币面额大于目标、哨兵溢出，并故意把容量方向写反构造最小反例。这样能确认循环方向来自题意，不是背模板。

\textbf{复杂度变化。} 暴力指数级；DP O(amount*coins.size()) 时间，O(amount) 空间。

\textbf{迁移追问。} 如果题目改成“零钱兑换 II 问组合数，循环语义不同”，先判断原状态是否仍然覆盖新答案集合，再决定是微调边界、改 value 语义，还是换算法。

\noindent\textbf{LeetCode 377 组合总和 IV。}

\textbf{题目说明。} LeetCode 377 组合总和 IV 的题面入口要先还原成具体任务：题目统计排列数，选择顺序不同算不同方案。

\textbf{样例入口。} 拿 nums=[1, 2, 3]、target=4 手算，1+3 和 3+1 是不同排列，所以 dp[4] 要按最后一步来自 1、2、3 分别累加。

\textbf{暴力基线。} 慢版本枚举每个物品选不选、选几次，并记录阶段和剩余容量。相同阶段和容量反复出现时，就可以把指数搜索压成表。放到本题就是：先按“题目统计排列数，选择顺序不同算不同方案”完整试慢版本；样例里已经能看到“规律是组合总和 IV 计数排列，容量在外层、数字在内层，循环顺序不能写成组合题”，所以优化前必须先能说清慢版本枚举了哪些候选。

\textbf{暴力过程。} 拿 nums=[1, 2, 3]、target=4 手算，1+3 和 3+1 是不同排列，所以 dp[4] 要按最后一步来自 1、2、3 分别累加。

\textbf{重复工作。} 题面模型：题目统计排列数，选择顺序不同算不同方案。暴力版的慢点要落到本题：慢版本会围绕“题目统计排列数，选择顺序不同算不同方案”枚举物品使用情况。样例规律是“规律是组合总和 IV 计数排列，容量在外层、数字在内层，循环顺序不能写成组合题”，说明阶段、容量和使用次数已经足够描述后续选择，没必要反复展开同一剩余容量。后面的优化只消掉这类重复工作，不能改变答案集合。

\textbf{优化条件。} 本题的关键观察是：外层遍历容量，内层遍历数字，让每个容量的方案按最后一个数累加。这句话必须能翻译成可维护状态；如果题目条件改变后观察不再成立，就不能继续套原来的写法。

\textbf{相邻状态。} 规律是组合总和 IV 计数排列，容量在外层、数字在内层，循环顺序不能写成组合题。把这个特例继续拆开看：阶段是物品，容量是资源，状态值是可达、最值还是计数。一个物品能否在同一轮再次使用，直接决定容量循环方向；组合和排列也由循环顺序表达。

\textbf{状态复用。} 把选择树压成阶段和容量；物品是否可重复使用决定一维表的遍历方向。本题落点是：规律是组合总和 IV 计数排列，容量在外层、数字在内层，循环顺序不能写成组合题。手算时只改被新输入影响的那一小部分，而不是回头重算完整候选。

\textbf{指针和字段。} 状态字段要能说清：物品阶段、容量、状态值、循环方向、取模或不可达哨兵；组合和排列必须用不同循环顺序表达。手算时按这个协议跟踪：拿 nums=[1, 2, 3]、target=4 手算，1+3 和 3+1 是不同排列，所以 dp[4] 要按最后一步来自 1、2、3 分别累加。然后按协议复查：手算时用一个容量很小的例子比较正序和倒序。只要循环方向改变后同一物品能否重复使用发生变化，就说明方向不是细节。

\textbf{代码落点。} 0-1 倒序、完全正序、排列外层容量、组合外层物品；每种循环方向都要能用一个小例子解释。关键代码行必须能逐句翻译成“旧状态怎样变成新状态”，否则就只是模板。

\textbf{正确性。} 证明从当前物品使用次数开始：0-1、完全、计数和最值的循环语义必须和题意一致。

\textbf{边界实验。} 先用target 为零、数字大于目标、方案数溢出、循环顺序做反例检查。实验时覆盖target 为零、数字大于目标、方案数溢出、循环顺序，并故意把容量方向写反构造最小反例。这样能确认循环方向来自题意，不是背模板。

\textbf{复杂度变化。} 暴力枚举物品选择指数级；背包 DP 通常按物品数乘容量或多维容量乘积计算，这是伪多项式复杂度。

\textbf{迁移追问。} 如果题目改成“若顺序不同不算新方案，就变成完全背包组合数，循环顺序要交换”，先判断原状态是否仍然覆盖新答案集合，再决定是微调边界、改 value 语义，还是换算法。

\noindent\textbf{LeetCode 416 分割等和子集。}

\textbf{题目说明。} LeetCode 416 分割等和子集 的题面入口要先还原成具体任务：总和一半是背包容量，每个数只能使用一次。

\textbf{样例入口。} 拿 [1, 5, 11, 5] 手算，总和 22，目标容量 11；可以用 11 单独凑成，也可以 5+5+1。

\textbf{暴力基线。} 暴力枚举每个数放到左集合还是右集合，最后检查两边和是否相等。

\textbf{暴力过程。} 以 [1, 5, 11, 5] 为例，总和 22，目标是找一个子集和为 11；可以选单独的 11，也可以选 5+5+1。

\textbf{重复工作。} 题面模型：总和一半是背包容量，每个数只能使用一次。暴力版的慢点要落到本题：浪费在不同选择顺序会反复得到相同剩余容量。问题只关心能否凑到 sum/2，不关心具体左右集合顺序。后面的优化只消掉这类重复工作，不能改变答案集合。

\textbf{优化条件。} 本题的关键观察是：0-1 背包可达性，一维容量必须倒序。这句话必须能翻译成可维护状态；如果题目条件改变后观察不再成立，就不能继续套原来的写法。

\textbf{相邻状态。} 规律是分割等和转成 0-1 背包是否能凑到 sum/2，容量必须倒序避免重复使用同一元素。把这个特例继续拆开看：阶段是物品，容量是资源，状态值是可达、最值还是计数。一个物品能否在同一轮再次使用，直接决定容量循环方向；组合和排列也由循环顺序表达。

\textbf{状态复用。} 先判断总和是否为偶数。dp[cap] 表示处理过若干数字后是否能凑出 cap，每个 num 只能用一次，所以容量倒序更新。手算时只改被新输入影响的那一小部分，而不是回头重算完整候选。

\textbf{指针和字段。} 状态字段要能说清：target=sum/2，cap 是容量，num 是当前物品；dp[cap-num] 为真则 dp[cap] 可变真。手算时按这个协议跟踪：处理 1 后可达 1；处理 5 后可达 5 和 6；处理 11 后 target 11 直接可达。

\textbf{代码落点。} 关键是 for cap from target down to num。若正序更新，会在同一轮把当前 num 用多次，变成完全背包。关键代码行必须能逐句翻译成“旧状态怎样变成新状态”，否则就只是模板。

\textbf{正确性。} 每个数只有选或不选两种状态；倒序容量保证转移读取的是上一阶段结果，不会重复使用当前数。

\textbf{边界实验。} 先用总和为奇数、目标为零、数值较大、复杂度依赖 target做反例检查。实验时覆盖总和为奇数、目标为零、数值较大、复杂度依赖 target，并故意把容量方向写反构造最小反例。这样能确认循环方向来自题意，不是背模板。

\textbf{复杂度变化。} 暴力 O(\ensuremath{2^{n}})，0-1 背包 O(n*target) 时间、O(target) 空间。

\textbf{迁移追问。} 如果题目改成“负数会破坏普通容量模型，需要重新建模”，先判断原状态是否仍然覆盖新答案集合，再决定是微调边界、改 value 语义，还是换算法。

\noindent\textbf{LeetCode 474 一和零。}

\textbf{题目说明。} LeetCode 474 一和零 的题面入口要先还原成具体任务：每个字符串消耗零和一两种容量，每个字符串最多选一次。

\textbf{样例入口。} 拿 strs=[10, 0, 1]、m=1、n=1 手算，字符串 10 消耗一个 0 和一个 1，单独就得到长度 1；0 和 1 也能组合成长度 2。

\textbf{暴力基线。} 慢版本枚举每个物品选不选、选几次，并记录阶段和剩余容量。相同阶段和容量反复出现时，就可以把指数搜索压成表。放到本题就是：先按“每个字符串消耗零和一两种容量，每个字符串最多选一次”完整试慢版本；样例里已经能看到“规律是二维 0-1 背包，容量是 0 的数量和 1 的数量，两维都要倒序”，所以优化前必须先能说清慢版本枚举了哪些候选。

\textbf{暴力过程。} 拿 strs=[10, 0, 1]、m=1、n=1 手算，字符串 10 消耗一个 0 和一个 1，单独就得到长度 1；0 和 1 也能组合成长度 2。

\textbf{重复工作。} 题面模型：每个字符串消耗零和一两种容量，每个字符串最多选一次。暴力版的慢点要落到本题：慢版本会围绕“每个字符串消耗零和一两种容量，每个字符串最多选一次”枚举物品使用情况。样例规律是“规律是二维 0-1 背包，容量是 0 的数量和 1 的数量，两维都要倒序”，说明阶段、容量和使用次数已经足够描述后续选择，没必要反复展开同一剩余容量。后面的优化只消掉这类重复工作，不能改变答案集合。

\textbf{优化条件。} 本题的关键观察是：二维 0-1 背包，两个容量都要倒序遍历。这句话必须能翻译成可维护状态；如果题目条件改变后观察不再成立，就不能继续套原来的写法。

\textbf{相邻状态。} 规律是二维 0-1 背包，容量是 0 的数量和 1 的数量，两维都要倒序。把这个特例继续拆开看：阶段是物品，容量是资源，状态值是可达、最值还是计数。一个物品能否在同一轮再次使用，直接决定容量循环方向；组合和排列也由循环顺序表达。

\textbf{状态复用。} 把选择树压成阶段和容量；物品是否可重复使用决定一维表的遍历方向。本题落点是：规律是二维 0-1 背包，容量是 0 的数量和 1 的数量，两维都要倒序。手算时只改被新输入影响的那一小部分，而不是回头重算完整候选。

\textbf{指针和字段。} 状态字段要能说清：物品阶段、容量、状态值、循环方向、取模或不可达哨兵；组合和排列必须用不同循环顺序表达。手算时按这个协议跟踪：拿 strs=[10, 0, 1]、m=1、n=1 手算，字符串 10 消耗一个 0 和一个 1，单独就得到长度 1；0 和 1 也能组合成长度 2。然后按协议复查：手算时用一个容量很小的例子比较正序和倒序。只要循环方向改变后同一物品能否重复使用发生变化，就说明方向不是细节。

\textbf{代码落点。} 0-1 倒序、完全正序、排列外层容量、组合外层物品；每种循环方向都要能用一个小例子解释。关键代码行必须能逐句翻译成“旧状态怎样变成新状态”，否则就只是模板。

\textbf{正确性。} 证明从当前物品使用次数开始：0-1、完全、计数和最值的循环语义必须和题意一致。

\textbf{边界实验。} 先用字符串全零或全一、容量为零、倒序方向、计数预处理做反例检查。实验时覆盖字符串全零或全一、容量为零、倒序方向、计数预处理，并故意把容量方向写反构造最小反例。这样能确认循环方向来自题意，不是背模板。

\textbf{复杂度变化。} 暴力枚举物品选择指数级；背包 DP 通常按物品数乘容量或多维容量乘积计算，这是伪多项式复杂度。

\textbf{迁移追问。} 如果题目改成“这是背包从一维容量扩展到二维容量的代表”，先判断原状态是否仍然覆盖新答案集合，再决定是微调边界、改 value 语义，还是换算法。

\noindent\textbf{LeetCode 494 目标和。}

\textbf{题目说明。} LeetCode 494 目标和 的题面入口要先还原成具体任务：正负号选择可代数转换为选若干数凑出特定正和。

\textbf{样例入口。} 拿 [1, 1, 1, 1, 1]、target=3 手算，给其中一个 1 加负号，其余为正可以得到 3，共有 5 种位置选择。也可转成选正数子集和为 (sum+target)/2。

\textbf{暴力基线。} 暴力递归给每个数添加 + 或 -，枚举 \ensuremath{2^{n}} 个符号序列并统计和为 target 的方案。

\textbf{暴力过程。} 以 [1, 1, 1, 1, 1]、target=3 为例，只有一个 1 取负、其余取正时可得到 3，共 5 种位置选择。

\textbf{重复工作。} 题面模型：正负号选择可代数转换为选若干数凑出特定正和。暴力版的慢点要落到本题：浪费在符号序列会反复得到相同部分和。把正数集合 P、负数集合 N 写成 sum(P)-sum(N)=target，可转成子集和计数。后面的优化只消掉这类重复工作，不能改变答案集合。

\textbf{优化条件。} 本题的关键观察是：若 S 加 target 为奇数或目标绝对值过大，直接无解。这句话必须能翻译成可维护状态；如果题目条件改变后观察不再成立，就不能继续套原来的写法。

\textbf{相邻状态。} 规律是符号分配变成 0-1 计数背包，奇偶和范围先判定。把这个特例继续拆开看：阶段是物品，容量是资源，状态值是可达、最值还是计数。一个物品能否在同一轮再次使用，直接决定容量循环方向；组合和排列也由循环顺序表达。

\textbf{状态复用。} 总和 total。若 total+target 为奇数或 abs(target)>total，直接 0；否则找子集和 positive=(total+target)/2 的方案数。手算时只改被新输入影响的那一小部分，而不是回头重算完整候选。

\textbf{指针和字段。} 状态字段要能说清：dp[cap] 是凑出 cap 的方案数，num 是当前数；0-1 计数背包容量倒序更新。手算时按这个协议跟踪：target=3、total=5，positive=4；选择四个 1 作为正数等价于选择哪个 1 不放进去，共 5 种。

\textbf{代码落点。} 关键是 dp[0]=1，for cap from positive down to num。数字 0 会让方案数翻倍，倒序更新能自然处理。关键代码行必须能逐句翻译成“旧状态怎样变成新状态”，否则就只是模板。

\textbf{正确性。} 每个符号分配和一个正数子集一一对应；子集和满足 positive 时，正负差正好为 target。

\textbf{边界实验。} 先用零元素会让方案数翻倍、目标为负、容量为零做反例检查。实验时覆盖零元素会让方案数翻倍、目标为负、容量为零，并故意把容量方向写反构造最小反例。这样能确认循环方向来自题意，不是背模板。

\textbf{复杂度变化。} 暴力 O(\ensuremath{2^{n}})，背包 O(n*positive) 时间，O(positive) 空间。

\textbf{迁移追问。} 如果题目改成“也可用哈希 DP 保存所有可能和，但背包转换更高效”，先判断原状态是否仍然覆盖新答案集合，再决定是微调边界、改 value 语义，还是换算法。

\noindent\textbf{LeetCode 518 零钱兑换 II。}

\textbf{题目说明。} LeetCode 518 零钱兑换 II 的题面入口要先还原成具体任务：硬币无限使用，问组合数而不是最少数量。

\textbf{样例入口。} 拿 amount=5、coins=[1, 2, 5] 手算，组合 5、2+2+1、2+1+1+1、全 1 一共 4 种；1+2 和 2+1 不能重复算。

\textbf{暴力基线。} 暴力 DFS 枚举每种硬币使用几个，或者枚举所有硬币排列再去重。

\textbf{暴力过程。} amount=5、coins=[1, 2, 5] 时，组合有 5、2+2+1、2+1+1+1、1+1+1+1+1，共 4 种；1+2 和 2+1 不能重复算。

\textbf{重复工作。} 题面模型：硬币无限使用，问组合数而不是最少数量。暴力版的慢点要落到本题：浪费在排列顺序重复。组合计数应固定硬币顺序，让每种组合只在一种顺序下被统计。后面的优化只消掉这类重复工作，不能改变答案集合。

\textbf{优化条件。} 本题的关键观察是：外层遍历硬币、内层金额正序，保证组合按固定硬币顺序统计。这句话必须能翻译成可维护状态；如果题目条件改变后观察不再成立，就不能继续套原来的写法。

\textbf{相邻状态。} 规律是组合计数时外层枚举硬币，内层容量正序。把这个特例继续拆开看：阶段是物品，容量是资源，状态值是可达、最值还是计数。一个物品能否在同一轮再次使用，直接决定容量循环方向；组合和排列也由循环顺序表达。

\textbf{状态复用。} 完全背包计数。dp[amount] 表示当前已处理硬币集合能凑出该金额的组合数；外层枚举 coin，内层金额正序。手算时只改被新输入影响的那一小部分，而不是回头重算完整候选。

\textbf{指针和字段。} 状态字段要能说清：coin 是当前允许使用的硬币，cap 是当前金额，dp[cap-coin] 是加入一枚 coin 前的组合数。手算时按这个协议跟踪：处理 coin=2 时，dp[4] 会从 dp[2] 来，表示使用两个 2；不会再把 1+2 和 2+1 分开统计。

\textbf{代码落点。} 关键循环是 for coin in coins: for cap from coin to amount: dp[cap]+=dp[cap-coin]。若外层金额内层硬币，会统计排列数。关键代码行必须能逐句翻译成“旧状态怎样变成新状态”，否则就只是模板。

\textbf{正确性。} 外层固定硬币顺序后，每个组合按最大硬币处理阶段唯一计入；金额正序允许当前硬币重复使用。

\textbf{边界实验。} 先用amount 为零、无硬币、组合数较大、循环顺序反例做反例检查。实验时覆盖amount 为零、无硬币、组合数较大、循环顺序反例，并故意把容量方向写反构造最小反例。这样能确认循环方向来自题意，不是背模板。

\textbf{复杂度变化。} 暴力指数级；完全背包计数 O(amount*coins.size()) 时间，O(amount) 空间。

\textbf{迁移追问。} 如果题目改成“若外层金额内层硬币，会统计排列数”，先判断原状态是否仍然覆盖新答案集合，再决定是微调边界、改 value 语义，还是换算法。

\noindent\textbf{LeetCode 879 盈利计划。}

\textbf{题目说明。} LeetCode 879 盈利计划 的题面入口要先还原成具体任务：每个项目消耗人数并贡献利润，利润达到下限后可截断。

\textbf{样例入口。} 拿 n=5、minProfit=3、group=[2, 2]、profit=[2, 3] 手算，第二个项目单独就达标，两个项目一起也达标且人数不超。

\textbf{暴力基线。} 慢版本枚举每个物品选不选、选几次，并记录阶段和剩余容量。相同阶段和容量反复出现时，就可以把指数搜索压成表。放到本题就是：先按“每个项目消耗人数并贡献利润，利润达到下限后可截断”完整试慢版本；样例里已经能看到“规律是状态包含已用人数和被截断到 minProfit 的利润，项目只能选一次”，所以优化前必须先能说清慢版本枚举了哪些候选。

\textbf{暴力过程。} 拿 n=5、minProfit=3、group=[2, 2]、profit=[2, 3] 手算，第二个项目单独就达标，两个项目一起也达标且人数不超。

\textbf{重复工作。} 题面模型：每个项目消耗人数并贡献利润，利润达到下限后可截断。暴力版的慢点要落到本题：慢版本会围绕“每个项目消耗人数并贡献利润，利润达到下限后可截断”枚举物品使用情况。样例规律是“规律是状态包含已用人数和被截断到 minProfit 的利润，项目只能选一次”，说明阶段、容量和使用次数已经足够描述后续选择，没必要反复展开同一剩余容量。后面的优化只消掉这类重复工作，不能改变答案集合。

\textbf{优化条件。} 本题的关键观察是：二维 DP 用人数和已达到利润表示方案数，利润维度超过目标就压到目标。这句话必须能翻译成可维护状态；如果题目条件改变后观察不再成立，就不能继续套原来的写法。

\textbf{相邻状态。} 规律是状态包含已用人数和被截断到 minProfit 的利润，项目只能选一次。把这个特例继续拆开看：阶段是物品，容量是资源，状态值是可达、最值还是计数。一个物品能否在同一轮再次使用，直接决定容量循环方向；组合和排列也由循环顺序表达。

\textbf{状态复用。} 把选择树压成阶段和容量；物品是否可重复使用决定一维表的遍历方向。本题落点是：规律是状态包含已用人数和被截断到 minProfit 的利润，项目只能选一次。手算时只改被新输入影响的那一小部分，而不是回头重算完整候选。

\textbf{指针和字段。} 状态字段要能说清：物品阶段、容量、状态值、循环方向、取模或不可达哨兵；组合和排列必须用不同循环顺序表达。手算时按这个协议跟踪：拿 n=5、minProfit=3、group=[2, 2]、profit=[2, 3] 手算，第二个项目单独就达标，两个项目一起也达标且人数不超。然后按协议复查：手算时用一个容量很小的例子比较正序和倒序。只要循环方向改变后同一物品能否重复使用发生变化，就说明方向不是细节。

\textbf{代码落点。} 0-1 倒序、完全正序、排列外层容量、组合外层物品；每种循环方向都要能用一个小例子解释。关键代码行必须能逐句翻译成“旧状态怎样变成新状态”，否则就只是模板。

\textbf{正确性。} 证明从当前物品使用次数开始：0-1、完全、计数和最值的循环语义必须和题意一致。

\textbf{边界实验。} 先用minProfit 为零、人数不足、方案数取模、项目只能选一次做反例检查。实验时覆盖minProfit 为零、人数不足、方案数取模、项目只能选一次，并故意把容量方向写反构造最小反例。这样能确认循环方向来自题意，不是背模板。

\textbf{复杂度变化。} 暴力枚举物品选择指数级；背包 DP 通常按物品数乘容量或多维容量乘积计算，这是伪多项式复杂度。

\textbf{迁移追问。} 如果题目改成“这是 0-1 背包的计数扩展，目标维度需要饱和处理”，先判断原状态是否仍然覆盖新答案集合，再决定是微调边界、改 value 语义，还是换算法。

\noindent\textbf{LeetCode 956 最高的广告牌。}

\textbf{题目说明。} LeetCode 956 最高的广告牌 的题面入口要先还原成具体任务：两组钢筋高度相等时得到广告牌，高度差是关键状态。

\textbf{样例入口。} 拿 rods=[1, 2, 3, 6] 手算，左塔用 6，右塔用 1+2+3，也能达到高度 6。

\textbf{暴力基线。} 慢版本枚举每个物品选不选、选几次，并记录阶段和剩余容量。相同阶段和容量反复出现时，就可以把指数搜索压成表。放到本题就是：先按“两组钢筋高度相等时得到广告牌，高度差是关键状态”完整试慢版本；样例里已经能看到“规律是状态保存两边高度差对应的较矮塔最大高度，加入一根杆可以放左、放右或不用”，所以优化前必须先能说清慢版本枚举了哪些候选。

\textbf{暴力过程。} 拿 rods=[1, 2, 3, 6] 手算，左塔用 6，右塔用 1+2+3，也能达到高度 6。

\textbf{重复工作。} 题面模型：两组钢筋高度相等时得到广告牌，高度差是关键状态。暴力版的慢点要落到本题：慢版本会围绕“两组钢筋高度相等时得到广告牌，高度差是关键状态”枚举物品使用情况。样例规律是“规律是状态保存两边高度差对应的较矮塔最大高度，加入一根杆可以放左、放右或不用”，说明阶段、容量和使用次数已经足够描述后续选择，没必要反复展开同一剩余容量。后面的优化只消掉这类重复工作，不能改变答案集合。

\textbf{优化条件。} 本题的关键观察是：DP 保存某个高度差下较矮一侧的最大高度，新增钢筋可放高侧、低侧或不用。这句话必须能翻译成可维护状态；如果题目条件改变后观察不再成立，就不能继续套原来的写法。

\textbf{相邻状态。} 规律是状态保存两边高度差对应的较矮塔最大高度，加入一根杆可以放左、放右或不用。把这个特例继续拆开看：阶段是物品，容量是资源，状态值是可达、最值还是计数。一个物品能否在同一轮再次使用，直接决定容量循环方向；组合和排列也由循环顺序表达。

\textbf{状态复用。} 把选择树压成阶段和容量；物品是否可重复使用决定一维表的遍历方向。本题落点是：规律是状态保存两边高度差对应的较矮塔最大高度，加入一根杆可以放左、放右或不用。手算时只改被新输入影响的那一小部分，而不是回头重算完整候选。

\textbf{指针和字段。} 状态字段要能说清：物品阶段、容量、状态值、循环方向、取模或不可达哨兵；组合和排列必须用不同循环顺序表达。手算时按这个协议跟踪：拿 rods=[1, 2, 3, 6] 手算，左塔用 6，右塔用 1+2+3，也能达到高度 6。然后按协议复查：手算时用一个容量很小的例子比较正序和倒序。只要循环方向改变后同一物品能否重复使用发生变化，就说明方向不是细节。

\textbf{代码落点。} 0-1 倒序、完全正序、排列外层容量、组合外层物品；每种循环方向都要能用一个小例子解释。关键代码行必须能逐句翻译成“旧状态怎样变成新状态”，否则就只是模板。

\textbf{正确性。} 证明从当前物品使用次数开始：0-1、完全、计数和最值的循环语义必须和题意一致。

\textbf{边界实验。} 先用空数组、无法平衡、重复长度、差值更新顺序做反例检查。实验时覆盖空数组、无法平衡、重复长度、差值更新顺序，并故意把容量方向写反构造最小反例。这样能确认循环方向来自题意，不是背模板。

\textbf{复杂度变化。} 暴力枚举物品选择指数级；背包 DP 通常按物品数乘容量或多维容量乘积计算，这是伪多项式复杂度。

\textbf{迁移追问。} 如果题目改成“这是背包从容量转成差值状态的代表”，先判断原状态是否仍然覆盖新答案集合，再决定是微调边界、改 value 语义，还是换算法。

\noindent\textbf{LeetCode 1049 最后一块石头的重量 II。}

\textbf{题目说明。} LeetCode 1049 最后一块石头的重量 II 的题面入口要先还原成具体任务：把石头分成两堆，最终差值等于两堆重量差。

\textbf{样例入口。} 拿 [2, 7, 4, 1, 8, 1] 手算，把石头分成两堆，最后差值越小结果越小；总和 23，尽量凑接近 11 的一堆。

\textbf{暴力基线。} 慢版本枚举每个物品选不选、选几次，并记录阶段和剩余容量。相同阶段和容量反复出现时，就可以把指数搜索压成表。放到本题就是：先按“把石头分成两堆，最终差值等于两堆重量差”完整试慢版本；样例里已经能看到“规律是 0-1 背包求不超过 sum/2 的最大可达重量”，所以优化前必须先能说清慢版本枚举了哪些候选。

\textbf{暴力过程。} 拿 [2, 7, 4, 1, 8, 1] 手算，把石头分成两堆，最后差值越小结果越小；总和 23，尽量凑接近 11 的一堆。

\textbf{重复工作。} 题面模型：把石头分成两堆，最终差值等于两堆重量差。暴力版的慢点要落到本题：慢版本会围绕“把石头分成两堆，最终差值等于两堆重量差”枚举物品使用情况。样例规律是“规律是 0-1 背包求不超过 sum/2 的最大可达重量”，说明阶段、容量和使用次数已经足够描述后续选择，没必要反复展开同一剩余容量。后面的优化只消掉这类重复工作，不能改变答案集合。

\textbf{优化条件。} 本题的关键观察是：0-1 背包找不超过总和一半的最大可达重量。这句话必须能翻译成可维护状态；如果题目条件改变后观察不再成立，就不能继续套原来的写法。

\textbf{相邻状态。} 规律是 0-1 背包求不超过 sum/2 的最大可达重量。把这个特例继续拆开看：阶段是物品，容量是资源，状态值是可达、最值还是计数。一个物品能否在同一轮再次使用，直接决定容量循环方向；组合和排列也由循环顺序表达。

\textbf{状态复用。} 把选择树压成阶段和容量；物品是否可重复使用决定一维表的遍历方向。本题落点是：规律是 0-1 背包求不超过 sum/2 的最大可达重量。手算时只改被新输入影响的那一小部分，而不是回头重算完整候选。

\textbf{指针和字段。} 状态字段要能说清：物品阶段、容量、状态值、循环方向、取模或不可达哨兵；组合和排列必须用不同循环顺序表达。手算时按这个协议跟踪：拿 [2, 7, 4, 1, 8, 1] 手算，把石头分成两堆，最后差值越小结果越小；总和 23，尽量凑接近 11 的一堆。然后按协议复查：手算时用一个容量很小的例子比较正序和倒序。只要循环方向改变后同一物品能否重复使用发生变化，就说明方向不是细节。

\textbf{代码落点。} 0-1 倒序、完全正序、排列外层容量、组合外层物品；每种循环方向都要能用一个小例子解释。关键代码行必须能逐句翻译成“旧状态怎样变成新状态”，否则就只是模板。

\textbf{正确性。} 证明从当前物品使用次数开始：0-1、完全、计数和最值的循环语义必须和题意一致。

\textbf{边界实验。} 先用单块石头、总和为奇偶、多个相同重量、容量方向做反例检查。实验时覆盖单块石头、总和为奇偶、多个相同重量、容量方向，并故意把容量方向写反构造最小反例。这样能确认循环方向来自题意，不是背模板。

\textbf{复杂度变化。} 暴力枚举物品选择指数级；背包 DP 通常按物品数乘容量或多维容量乘积计算，这是伪多项式复杂度。

\textbf{迁移追问。} 如果题目改成“它和分割等和子集同源，只是目标从是否可达变成最小差”，先判断原状态是否仍然覆盖新答案集合，再决定是微调边界、改 value 语义，还是换算法。

\noindent\textbf{LeetCode 1155 掷骰子等于目标和的方法数。}

\textbf{题目说明。} LeetCode 1155 掷骰子等于目标和的方法数 的题面入口要先还原成具体任务：多个骰子每个只能取一到 faces 的点数，问目标和方案数。

\textbf{样例入口。} 拿 2 个 6 面骰、target=7 手算，第一颗为 1 时第二颗必须 6，第一颗为 2 时第二颗必须 5。

\textbf{暴力基线。} 慢版本枚举每个物品选不选、选几次，并记录阶段和剩余容量。相同阶段和容量反复出现时，就可以把指数搜索压成表。放到本题就是：先按“多个骰子每个只能取一到 faces 的点数，问目标和方案数”完整试慢版本；样例里已经能看到“规律是阶段是骰子个数，容量是当前点数和，每个骰子枚举 1..k 的面值”，所以优化前必须先能说清慢版本枚举了哪些候选。

\textbf{暴力过程。} 拿 2 个 6 面骰、target=7 手算，第一颗为 1 时第二颗必须 6，第一颗为 2 时第二颗必须 5。

\textbf{重复工作。} 题面模型：多个骰子每个只能取一到 faces 的点数，问目标和方案数。暴力版的慢点要落到本题：慢版本会围绕“多个骰子每个只能取一到 faces 的点数，问目标和方案数”枚举物品使用情况。样例规律是“规律是阶段是骰子个数，容量是当前点数和，每个骰子枚举 1..k 的面值”，说明阶段、容量和使用次数已经足够描述后续选择，没必要反复展开同一剩余容量。后面的优化只消掉这类重复工作，不能改变答案集合。

\textbf{优化条件。} 本题的关键观察是：按骰子数量分阶段，容量为当前点数和，枚举当前骰子的点数转移。这句话必须能翻译成可维护状态；如果题目条件改变后观察不再成立，就不能继续套原来的写法。

\textbf{相邻状态。} 规律是阶段是骰子个数，容量是当前点数和，每个骰子枚举 1..k 的面值。把这个特例继续拆开看：阶段是物品，容量是资源，状态值是可达、最值还是计数。一个物品能否在同一轮再次使用，直接决定容量循环方向；组合和排列也由循环顺序表达。

\textbf{状态复用。} 把选择树压成阶段和容量；物品是否可重复使用决定一维表的遍历方向。本题落点是：规律是阶段是骰子个数，容量是当前点数和，每个骰子枚举 1..k 的面值。手算时只改被新输入影响的那一小部分，而不是回头重算完整候选。

\textbf{指针和字段。} 状态字段要能说清：物品阶段、容量、状态值、循环方向、取模或不可达哨兵；组合和排列必须用不同循环顺序表达。手算时按这个协议跟踪：拿 2 个 6 面骰、target=7 手算，第一颗为 1 时第二颗必须 6，第一颗为 2 时第二颗必须 5。然后按协议复查：手算时用一个容量很小的例子比较正序和倒序。只要循环方向改变后同一物品能否重复使用发生变化，就说明方向不是细节。

\textbf{代码落点。} 0-1 倒序、完全正序、排列外层容量、组合外层物品；每种循环方向都要能用一个小例子解释。关键代码行必须能逐句翻译成“旧状态怎样变成新状态”，否则就只是模板。

\textbf{正确性。} 证明从当前物品使用次数开始：0-1、完全、计数和最值的循环语义必须和题意一致。

\textbf{边界实验。} 先用target 太小或太大、取模、骰子数为一、滚动数组方向做反例检查。实验时覆盖target 太小或太大、取模、骰子数为一、滚动数组方向，并故意把容量方向写反构造最小反例。这样能确认循环方向来自题意，不是背模板。

\textbf{复杂度变化。} 暴力枚举物品选择指数级；背包 DP 通常按物品数乘容量或多维容量乘积计算，这是伪多项式复杂度。

\textbf{迁移追问。} 如果题目改成“这是有界选择计数，和完全背包不能混”，先判断原状态是否仍然覆盖新答案集合，再决定是微调边界、改 value 语义，还是换算法。

\noindent\textbf{LeetCode 1449 数位成本和为目标值的最大数字。}

\textbf{题目说明。} LeetCode 1449 数位成本和为目标值的最大数字 的题面入口要先还原成具体任务：每个数字可重复选择，容量是成本，目标是组成字典序和长度最大的数字。

\textbf{样例入口。} 拿 cost 中数字 7 和 2 成本较低、target=9 手算，先最大化能组成的位数，再从 9 到 1 尝试还原字典序最大的数字。

\textbf{暴力基线。} 慢版本枚举每个物品选不选、选几次，并记录阶段和剩余容量。相同阶段和容量反复出现时，就可以把指数搜索压成表。放到本题就是：先按“每个数字可重复选择，容量是成本，目标是组成字典序和长度最大的数字”完整试慢版本；样例里已经能看到“规律是完全背包求最大位数，重建时优先放大数字且保证剩余目标可达”，所以优化前必须先能说清慢版本枚举了哪些候选。

\textbf{暴力过程。} 拿 cost 中数字 7 和 2 成本较低、target=9 手算，先最大化能组成的位数，再从 9 到 1 尝试还原字典序最大的数字。

\textbf{重复工作。} 题面模型：每个数字可重复选择，容量是成本，目标是组成字典序和长度最大的数字。暴力版的慢点要落到本题：慢版本会围绕“每个数字可重复选择，容量是成本，目标是组成字典序和长度最大的数字”枚举物品使用情况。样例规律是“规律是完全背包求最大位数，重建时优先放大数字且保证剩余目标可达”，说明阶段、容量和使用次数已经足够描述后续选择，没必要反复展开同一剩余容量。后面的优化只消掉这类重复工作，不能改变答案集合。

\textbf{优化条件。} 本题的关键观察是：完全背包先最大化位数，再按数字从大到小贪心还原答案。这句话必须能翻译成可维护状态；如果题目条件改变后观察不再成立，就不能继续套原来的写法。

\textbf{相邻状态。} 规律是完全背包求最大位数，重建时优先放大数字且保证剩余目标可达。把这个特例继续拆开看：阶段是物品，容量是资源，状态值是可达、最值还是计数。一个物品能否在同一轮再次使用，直接决定容量循环方向；组合和排列也由循环顺序表达。

\textbf{状态复用。} 把选择树压成阶段和容量；物品是否可重复使用决定一维表的遍历方向。本题落点是：规律是完全背包求最大位数，重建时优先放大数字且保证剩余目标可达。手算时只改被新输入影响的那一小部分，而不是回头重算完整候选。

\textbf{指针和字段。} 状态字段要能说清：物品阶段、容量、状态值、循环方向、取模或不可达哨兵；组合和排列必须用不同循环顺序表达。手算时按这个协议跟踪：拿 cost 中数字 7 和 2 成本较低、target=9 手算，先最大化能组成的位数，再从 9 到 1 尝试还原字典序最大的数字。然后按协议复查：手算时用一个容量很小的例子比较正序和倒序。只要循环方向改变后同一物品能否重复使用发生变化，就说明方向不是细节。

\textbf{代码落点。} 0-1 倒序、完全正序、排列外层容量、组合外层物品；每种循环方向都要能用一个小例子解释。关键代码行必须能逐句翻译成“旧状态怎样变成新状态”，否则就只是模板。

\textbf{正确性。} 证明从当前物品使用次数开始：0-1、完全、计数和最值的循环语义必须和题意一致。

\textbf{边界实验。} 先用无法组成目标、多个数字同成本、target 为零、字符串比较做反例检查。实验时覆盖无法组成目标、多个数字同成本、target 为零、字符串比较，并故意把容量方向写反构造最小反例。这样能确认循环方向来自题意，不是背模板。

\textbf{复杂度变化。} 暴力枚举物品选择指数级；背包 DP 通常按物品数乘容量或多维容量乘积计算，这是伪多项式复杂度。

\textbf{迁移追问。} 如果题目改成“这题说明背包状态值不一定是数值收益，也可以是长度和字典序”，先判断原状态是否仍然覆盖新答案集合，再决定是微调边界、改 value 语义，还是换算法。

\subsection{实验复盘}

追问通常会问倒序和正序的区别、组合和排列的区别、为什么复杂度是伪多项式。请准备一个循环方向写反的最小反例。

复盘本节时先从 LeetCode 279 完全平方数、LeetCode 322 零钱兑换、LeetCode 377 组合总和 IV 开始：每道题都先手写一个能覆盖全部候选的暴力版，再指出优化结构具体保存了哪一段历史信息，最后用相邻案例比较为什么不能互换解法。
